\documentclass[11pt, a4paper]{article}
\usepackage{graphicx}
\usepackage{geometry}
\usepackage{biblatex}
\usepackage{parskip}
\usepackage{stmaryrd}
\usepackage{blindtext}
\usepackage{amsthm, amsmath, amssymb}
\usepackage{mathtools}
\usepackage{mathrsfs}
\usepackage{hyperref}
\usepackage{physics}
\usepackage{tikz}
\usepackage{nicematrix}
\usepackage{epigraph}
\usepackage{enumitem}

\geometry{margin=2.2cm}
\renewcommand{\qedsymbol}{$\blacksquare$}
\newcommand{\suchthat}{\;\ifnum\currentgrouptype=16 \middle\fi|\;}
\newcommand\multr{\stackrel{\mathclap{\normalfont\mbox{.}}}{\mathbb{R}}}

\makeatletter
\newcommand\longleftrightarrowfill@{%
  \arrowfill@\leftarrow\relbar\rightarrow}
\makeatother

% 1. Macro for bracketed matrix [] -> \bmat{...}
\newcommand{\bmat}[1]{\begin{bmatrix} #1 \end{bmatrix}}

% 2. Macro for single vertical bar matrix || -> \vmat{...}
\newcommand{\vmat}[1]{\begin{vmatrix} #1 \end{vmatrix}}

\newcommand{\RR}{\mathbb{R}}
\newcommand{\CC}{\mathbb{C}}
\newcommand{\ZZ}{\mathbb{Z}}

\newcommand{\adj}{\operatorname{adj}}
\newcommand{\diag}{\operatorname{diag}}
\newcommand{\inv}{^{-1}}
% \newcommand{\rank}{\operatorname{rank}}
\newcommand{\matmn}{M_{m \times n}(\RR)}
% \newcommand{\ker}{\operatorname{ker}}
\newcommand{\im}{\operatorname{im}}
\newcommand{\spn}{\operatorname{span}}
\newcommand{\id}{\operatorname{id}}
\newcommand{\tpse}{^{\mathsf{T}}}

\newcommand{\ul}{\underline}
\newcommand{\lline}{\noindent\rule{\textwidth}{0.4pt}}

\theoremstyle{definition}
\newtheorem{definition}{Definition}[section]
\newtheorem{problem}{Problem}[subsection]
\newtheorem{lemma}{Lemma}[section]
\newtheorem{example}{Example}[section]
\newtheorem{exercs}{Exercise}
\newtheorem{theorem}{Theorem}[section]
\newtheorem{remark}{Remark}[section]
\newtheorem{corollary}{Corollary}[section]
\newtheorem{fact}{Fact}[section]
\newenvironment{amatrix}[1]{%
  \left[\begin{array}{@{}*{#1}{c}|c@{}}
}{%
  \end{array}\right]
}

\hypersetup{
    colorlinks=true,
    linkcolor=blue,
    filecolor=magenta,      
    urlcolor=blue,
    pdftitle={Math 40 Notes},
    }

% \setlength{\droptitle}{-1.0cm} % move title up
\title{\textbf{Math 40 Notes}}
\author{\texttt{elsoliva@up.edu.ph}}
\date{}

\begin{document}
\maketitle
\section{Matrices and Matrix Operations}
% \epigraph{Immer mit den einfachsten Beispielen anfangen.}{--- David Hilbert}

\begin{definition}
An $m \times n$ matrix is a rectangular array of real numbers arranged in $m$ horizontal rows and $n$ vertical columns.
\end{definition}

$$A = \left [
\begin{matrix}
a_{11} & \dots & a_{1n} \\
\vdots & \ddots & \vdots \\
a_{m1} & \dots & a_{mn}
\end{matrix} \right ] = [a_{ij}]_{m \times n}
$$

\begin{definition}
An $m \times n$ matrix where $m, n \in \mathbb{R}^+$ and $m = n$, wherein it has the same number of rows and columns is called a \textbf{square matrix}. A $1 \times n$ matrix is called a \textbf{row vector}. And an $n \times 1$ matrix is called a \textbf{column vector}.
\end{definition}

\begin{definition}
We will define the set of all $m \times n$ matrices with $\mathbb{R}$ entries as $M_{m \times n}(\mathbb{R})$. And we will denote the set of all square matrices with size $n$ as $M_n(\mathbb{R})$.
\end{definition}

\begin{example}
Let $A$ be a square matrix with size $3 \times 3$, and 
$$ A = [a_{ij}]_{3 \times 3} = 
    \left [ 
    \begin{matrix}
    1 & 1 & 0 \\
    2 & 0 & 1 \\
    3 & -1 & 2
    \end{matrix} 
\right ]$$
Find the value of the following entries: $a_{12}, a_{31}$, and $a_{33}$.

Recall that $a_{ij}$ is the $(i, j)$ entry of $A$ where $i$ represents the row and $j$ is the column. For the example above use the subscript $ij$ as a guide. $a_{12}$ is the $(1, 2)$ entry in the matrix, thus, $a_{12} = 1$. The following entries are left to the reader.

\end{example}

\begin{example}
Let $B$ be a matrix wherein $B = [b_{ij}]_{3 \times 2}$ where $b_{ij} = i + j$, draw the matrix and its entries.

$$B = \bmat{2 & 3 \\ 3 & 4 \\ 4 & 5}$$
The $b_{ij}$ entries are left as an exercise to the reader.
\end{example}

\begin{remark}\label{remark:1.1}
Two matrices $A = [a_{ij}]$ and $B = [b_{ij}]$ are said to be equal, denoted by $A = B$, if they have the same size and $a_{ij} = b_{ij}$ for all $i, j$. 
\end{remark} 

However, by \hyperref[remark:1.1]{Remark 1.1}, this does not imply that two matrices with the same size are equal. Let us illustrate an example. Suppose that there exists two matrices $C, D \in M_n(\RR)$.
$$C = 
\setlength{\extrarowheight}{1mm}
\begin{bNiceMatrix}
-4 & 6 \\
1 & 7
\CodeAfter
  \tikz \draw (2-1) circle (1.7mm) 
              (2-2) circle (1.7mm) ;
\end{bNiceMatrix}, \quad
D = 
\setlength{\extrarowheight}{1mm}
\begin{bNiceMatrix}
-4 & 6 \\
7 & 1
\CodeAfter
  \tikz \draw (2-1) circle (1.7mm) 
              (2-2) circle (1.7mm) ;
\end{bNiceMatrix}
$$

Notice that they have different $(2, 1)$ and $(2, 2)$ entries. Thus, $[c_{ij}] \neq [d_{ij}]$.

\begin{definition}
Let $A = [a_{ij}]_{n \times n}$, wherein 
$$A = \left [
\begin{matrix}
a_{11} & \dots & a_{1n} \\
\vdots & \ddots & \vdots \\
a_{m1} & \dots & a_{mn}
\end{matrix} \right ] \in M_n(\mathbb{R}),$$

1. $a_{ij}$ where $i = j$ are called the \underline{main diagonal entries}

2. $A$ is said to be:
\begin{enumerate}
    \item[a.] \textbf{diagonal} if all entries $a_{ij} = 0$ where $i \neq j$,
    \item[b.] \textbf{upper triangular} if all entries $a_{ij} = 0$ where $i > j$,
    \item[c.] \textbf{lower triangular} if all entries $a_{ij} = 0$ where $i < j$.
\end{enumerate}
\end{definition}

\begin{definition}
An $n \times n$ matrix is called an \textbf{identity matrix}, denoted by $I_n$ wherein all the diagonal entries are equal to 1, and the rest are 0. 
$$I_n = \bmat{1 & 0 & \cdots & 0 \\ 0 & 1 & \cdots & 0 \\ 0 & 0 & \ddots & 0 \\ 0 & 0 & \cdots & 1} \in M_n(\RR)$$
An $m \times n$ \textbf{zero matrix}, denoted by $0_{m \times n}$ is a matrix wherein all its entries are 0.
\end{definition}

\begin{remark} Listed below are some observations on the matrices we have discussed:

1. $I_n = [\delta_{ij}]_{n \times n}$ where it is represented by a piecewise function\footnote{Also known as the Kronecker delta function.} $\delta_{ij} = \begin{cases}
    1 & \text{if $i = j$} \\
    0 & \text{if $i \neq j$}
\end{cases}$

2. Suppose you have a square matrix of size $n$, $A \in M_n(\mathbb{R})$ is diagonal with diagonal entries $a_{11}, a_{22}, \dots, a_{nn}$, then we write
$$A = \diag(a_{11}, a_{22}, \dots, a_{nn})$$
\end{remark}

\lline

To illustrate examples of diagonal matrices:
$$P = \diag(2, \pi, e, 0) = \bmat{2 & 0 & 0 & 0 \\ 0 & \pi & 0 & 0 \\ 0 & 0 & e & 0 \\ 0 & 0 & 0 & 0}$$
$$I_2 = \diag(1, 1) = \bmat{1 & 0 \\ 0 & 1}$$
$$0_{3 \times 3} = \diag(0, 0, 0) = \bmat{0 & 0 & 0 \\ 0 & 0 & 0 \\ 0 & 0 & 0}$$

Do note that the notation for diagonal matrices \textbf{only applies to square matrices}. And in general, suppose there exists an identity matrix $I_n$ it can be represented with the diagonal notation with $\diag(\underbrace{1, 1, \dots, 1}_n)$ (this also applies to the zero matrix but replace $1$ with $0$).

\begin{definition}
Let us consider two matrices with $m \times n$ size $A = [a_{ij}]$ and $B = [b_{ij}]$, $A, B \in M_{m \times n}(\mathbb{R})$ and a scalar $c \in \mathbb{R}$

1. The sum of $A$ and $B$, 
$$A + B = [a_{ij} + b_{ij}]_{m \times n}$$

2. The scalar product of $c$ and $A$, 
$$cA = [ca_{ij}]_{m \times n}$$

3. The transpose of $A$,
$$A^{\mathsf{T}} = [t_{ij}]_{n \times m}$$
where $t_{ij} = a_{ji}$.
\end{definition}

\begin{example}
Let us illustrate the definitions we have shown above:
\begin{enumerate}
    \item \begin{align*}
        \bmat{4 & -3 \\ 6 & 0 \\ 7 & -1} + 3\bmat{1 & 1 \\ -2 & 5 \\ 0 & -2} &= \bmat{4 & -3 \\ 6 & 0 \\ 7 & -1} + \bmat{3 & 3 \\ -6 & 15 \\ 0 & -6} \\
        & = \bmat{7 & 0 \\ 0 & 15 \\ 7 & -7}
    \end{align*}
    \item $$\bmat{-1 & 0 \\ 4 & 1 \\ 2 & -3}^{\mathsf{T}} = \bmat{-1 & 4 & 2 \\ 0 & 1 & -3}$$
    \item $$\bmat{-1 & 2 & 0 \\ 2 & 1 & 0 \\ 0 &  0 & 3}\tpse = \bmat{-1 & 2 & 0 \\ 2 & 1 & 0 \\ 0 & 0 & 3}$$
    Notice that in this example the transpose of a matrix is equal to itself. We call such matrices \textbf{symmetric}, we will further expand on this on \hyperref[def:14.4]{Chapter 14}.
    \item $$\frac{1}{2}\bmat{4 & 2 & 0}\tpse - \bmat{1 \\ -4 \\ 2} = \bmat{2 \\ 1 \\ 0} - \bmat{1 \\ -4 \\ 2} = \bmat{1 \\ 5 \\ -2}$$
\end{enumerate}
\end{example}

\begin{theorem}[Properties of Matrix Addition, Scalar Multiplication, and Transposition]
\leavevmode \\
Let $A, B, C \in M_{m \times n}(\mathbb{R})$ and $c, d \in \mathbb{R}$. Then,

1. Matrix addition is commutative and associative,
$$A + B = B + A$$
$$A + (B + C) = (A + B) + C$$

2. There exists an identity matrix $0_{m \times n}$ such that when added to $A$ it returns $A$,
$$A + 0_{m \times n} = A$$

3. There exists an additive inverse $-A = -1(A)$, such that when added to $A$ it returns $0_{m \times n}$,
$$A + (-A) = 0_{m \times n}$$

4. Scalars could commute,
$$(cd)A = c(dA) = d(cA)$$

5. Scalars could be distributed,
$$(c + d)A = cA + dA$$
$$c(A + B) = cA + cB$$

6. $$(A + B)^{\mathsf{T}} = A^{\mathsf{T}} + B^{\mathsf{T}}$$

7. $$(cA)^{\mathsf{T}} = cA^{\mathsf{T}}$$

8. $$(A^{\mathsf{T}})^{\mathsf{T}} = A$$
\end{theorem}

\begin{example}
Let $A = \bmat{-1 & 1 \\ 2 & -3 \\ 4 & -5}$ and $B = \bmat{0 & 1 \\ 5 & 4 \\ 3 & -2}$. Solve for $X$ in the matrix equation $4X\tpse - 3A = B$.
\begin{align*}
4X\tpse = B + 3A &\Longrightarrow X\tpse = \frac{1}{4}(B + 3A) \\
&\Longrightarrow (X\tpse)\tpse = \left[\frac{1}{4}(B + 3A)\right]\tpse \\
X = \left[\frac{1}{4}B + \frac{3}{4}A\right]\tpse &= \left[\bmat{0 & 1/4 \\ 5/4 & 1 \\ 3/4 & -1/2} + \bmat{-3/4 & 3/4 \\ 3/2 & -9/4 \\ 3 & -15/4}\right]\tpse \\\\
&=\bmat{-3/4 & 11/4 & 15/4 \\ 1 & -5/4 & -17/4}
\end{align*}
\end{example}


\begin{definition}
Let $A = [a_{ij}]_{m \times n}$ and $B = [b_{ij}]_{n \times p}$, we define a product of $A$ and $B$, denoted by $A \times B = [c_{ij}]_{m \times p}$  where $\displaystyle c_{ij} = \bmat{a_{i1} & a_{i2} & \cdots a_{in}} \bmat{b_{1j} \\ \vdots \\ b_{nj}} = a_{i1}b_{1j} + a_{i2}b_{2j} + \cdots + a_{in}b_{nj}$.
\end{definition}

\begin{example}
Find $AB$ and $BA$\footnote{The underbraces under the matrices are simply visuals and should not be placed when solving.}
\begin{enumerate}
    \item $$A = \underbrace{\bmat{1 & 2 & -1 \\ 3 & 1 & 4}}_{2 \times 3}, \quad B = \underbrace{\bmat{-2 & 5 \\ 4 & -3 \\ 2 & 1}}_{3 \times 2}$$
    \begin{enumerate}
        \item \begin{align*}
            AB &= \bmat{1(-2) + 2(4) + (-1)(2) & 1(5) + 2(-3) + (-1)(1) \\ 3(-2) + 1(4) + (-1)(2) & 3(5) + 1(-3) + 4(1)} \\
            &= \bmat{4 & -2 \\ 6 & 16}
        \end{align*}

        \item \begin{align*}
            BA &= \bmat{-2 & 5 \\ 4 & -3 \\ 2 & 1} \bmat{1 & 2 & -1 \\ 3 & 1 & 4} \\
            &= \bmat{-2(1) + 5(3) & -2(2) + 5(1) & -2(-1) + 5(4) \\ 4(1) + (-3)(3) & 4(2) + (-3)(1) & 4(-1) + (-3)(4) \\ 2(1) + 1(3) & 2(2) + 1(1) & 2(-1) + 1(4)} \\
            & = \bmat{13 & 1 & 22 \\ -5 & 5 & -16 \\ 5 & 5 & 2}
        \end{align*}
    \end{enumerate}

    \item $$A = \bmat{1 & 2 \\ 1 & 1}, \quad B = \bmat{-1 & 2 \\ 1 & -1}$$
    \begin{enumerate}
        \item \begin{align*}
            AB &= \bmat{1(-1) + 2(1) & 1(2) + 2(-1) \\ 1(-1) + 1(1) & 1(2) + 1(-1)} \\
            &=\bmat{1 & 0 \\ 0 & 1}
        \end{align*}
        \item \begin{align*}
            BA &= \bmat{-1(1) + 2(1) & -1(2) + 2(1) \\ 1(1) + (-1)(1) & 1(2) + (-1)(1)} \\
            &= \bmat{1 & 0 \\ 0 & 1}
        \end{align*}
    \end{enumerate}
\end{enumerate}
\end{example}

\begin{theorem}
Let $A, B, C$ be matrices of suitable size and $c \in \mathbb{R}$,

\begin{enumerate}
    \item If $A \in M_{m \times n}(\mathbb{R})$ then
    \begin{enumerate}
        \item $I_mA = A = AI_n$
        \item $A0_{n \times p} =0_{m \times p}$ and $0_{q \times m}A = 0_{q \times n}$
    \end{enumerate}
    \item $(AB)C = A(BC)$
    \item $A(B + C) = AB + AC$ and $(A + B)C = AC + BC$
    \item $c(AB) = (cA)B = A(cB)$
    \item $(AB)\tpse = B\tpse A\tpse$
\end{enumerate}
\end{theorem}
\newpage
\section{Solving Systems of Linear Equations}
Consider a system of linear equations with $n$ unknowns and $m$ equations:
$$(*) = 
\begin{cases}
\begin{aligned}
a_{11}x_1 & + a_{12}x_2 + \cdots + a_{1n}x_n && = b_1 \\
a_{21}x_1 & + a_{22}x_2 + \cdots + a_{2n}x_n && = b_2 \\
&\ \vdots && \ \vdots \\
a_{m1}x_1 & + a_{m2}x_2 + \cdots + a_{mn}x_n && = b_m
\end{aligned}
\end{cases}
$$

\begin{theorem}
For $(*)$ exactly one of the following holds:
\begin{itemize}
    \item $(*)$ has exactly one solution (consistent and independent)
    \item $(*)$ has infinitely many solutions (consistent and dependent)
    \item $(*)$ has no solutions (inconsistent)
\end{itemize}
\end{theorem}

\begin{example}
We could rewrite $(*)$ into its matrix form, $Ax = b$, where $A$ is the coefficient matrix and $b$ is the vector of constants,
$$\underbrace{\bmat{a_{11} & a_{12} & \cdots & a_{1n} \\ a_{21} & a_{22} & \cdots & a_{2n} \\ \vdots & \vdots & \ddots & \vdots \\ a_{m1} & a_{m2} & \cdots & a_{mn}}}_A \underbrace{\bmat{x_1 \\ x_2 \\ \vdots \\ x_n}}_x = \underbrace{\bmat{b_1 \\ b_2 \\ \vdots \\ b_n}}_b$$
and we can denote the augmented matrix $[A \mid b]$ as
$$\begin{amatrix}{4}
a_{11} & a_{12} & \cdots & a_{1n} & b_1 \\
a_{21} & a_{22} & \cdots & a_{2n} & b_2 \\
\vdots & \vdots & \ddots & \vdots & \vdots \\
a_{m1} & a_{m2} & \cdots & a_{mn} & b_n \\
\end{amatrix}
$$
Notice that we have placed the coefficients in the left side and is separated by a vertical bar.\footnote{This serves as a visual representation on what we're separating.} On the right side we have the constants.
\end{example}

\begin{example}
Consider the following system of equations
$$\begin{cases}
    w + 3x + y + z = 3 \\
    2w + 2x + y + 2z = 8 \\
    3w + x + 2y - z = -1
\end{cases}$$
\end{example}

We can rewrite this in matrix form as:
$$\left[ 
\begin{matrix}
1 & 3 & 1 & 1 \\
2 & 2 & 1 & 2 \\
3 & 1 & 2 & -1
\end{matrix}
\right]
\left[
\begin{matrix}
w \\
x \\
y \\
z \\
\end{matrix}
\right] = 
\left[
\begin{matrix}
3 \\
8 \\
-1
\end{matrix}
\right]
$$

Shown below is its form in augmented matrix:
$$
\begin{amatrix}{4}
1 & 3 & 1 & 1 & 3 \\
3 & 1 & 2 & -1 & -1 \\
\end{amatrix}$$

\newpage
\underline{\textbf{Recall}}

In your previous math classes (Pre-calculus), you have learned solving system of equations through elimination by addition. To solve a system of equations, the goal is to convert it to an easier form that is easy to solve.

\underline{\textbf{Question}.} Which processes will yield equivalent system of equations?

Suppose we have the following system of equations:
$$\begin{cases}
    x - y = -2 \\
    2x + 6y = 12
\end{cases}$$

I. Switching the position of the equations:
\begin{align*}
\begin{cases}
    x - y = -2 \\
    2x + 6y = 12
\end{cases} \longleftrightarrow\,\,\,\,\,
\begin{cases}
    2x + 6y = 12 \\
    x - y = -2
\end{cases}
\end{align*}

II. Multiplying by a nonzero constant,
\begin{align*}
\begin{cases}
    x - y = -2 \\
    2x + 6y = 12
\end{cases} \longleftrightarrow \,\,\,\,\,
\begin{cases}
    2x + 6y = 12 \\
    x - y = -2
\end{cases}
\end{align*}


III. Add a constant multiple of one equation to another
\begin{align*}
\begin{cases}
    x - y = -2 \\
    2x + 6y = 12
\end{cases} 
&\longleftrightarrow \,\,\,\,\,
\begin{cases}
    x - y = -2 \\
    8y = 16
\end{cases} \\[1em]
&\longleftrightarrow \,\,\,\,\,
\begin{cases}
    x - y = -2 \\
    y = 2
\end{cases}
\end{align*}
All of these operations will yield the equivalent system of equations!

Suppose we take the system of equations from above and convert it to an augmented matrix,
$$\begin{amatrix}{2}
1 & 1 & - 2\\ 0 & 1 & 2 \\
\end{amatrix}
$$
Notice that it has a staircase form shown, as 
$$
\begin{bNiceMatrix}
1 & 1 & -2 \\ 
0 & 1 & \hspace{0.3cm}2 
\CodeAfter
  \tikz \draw[red, semithick, fill=red, fill opacity=0.1] 
    (1-1.north west) -- (1-3.north east) -- 
    (1-3.south east) -- (2-3.south east) -- 
    (2-2.south west) -- (1-2.south west) -- 
    (1-1.south west) -- cycle;
\end{bNiceMatrix}
$$
This form is also known as the \href{https://math.stackexchange.com/questions/1628192/what-does-echelon-mean/1629207#1629207}{row echelon form}.
\vspace{1em}

\begin{definition}
An $m \times n$ matrix is said to be in \textbf{row echelon form} if the following holds:
\begin{itemize}
    \item All zero rows of $A$ (if any) are at the bottom of $A$;
    \item The leftmost nonzero entry of every nonzero row of $A$ is 1;
    \item The leading entry of a higher row is strictly to the left of the leading entry of the next row.
\end{itemize}
\end{definition}

\begin{remark}
$A \in M_{m \times n}(\mathbb{R)}$ is in RREF (reduced row echelon form) if $A$ is in REF (reduced echelon form) and every leading entry of $A$ is the only nonzero entry of its column.
\end{remark}

\newpage
\begin{remark}
Each equation in a system corresponds to a row in its augmented matrix. The following processes will now yield equivalent systems:
\begin{itemize}
    \item \textbf{Type I} - Switch two rows in augmented matrix;
    \item \textbf{Type II} - Multiply a row in augmented matrix by nonzero constant;
    \item \textbf{Type III} - Add a constant multiple of one row in the augmented matrix to another row.
\end{itemize}
\end{remark}

\subsection{Gauss-Jordan Reduction}
To solve the system $Ax = b$:
\begin{itemize}
    \item Form the augmented matrix $[A \mid b]$
    \item Use elementary row operations to bring augmented matrix to RREF
    \item Express the resulting matrix as a system of linear equations
    \item Use back-substitution
\end{itemize}

\begin{remark}
\textbf{Gaussian elimination} is the term when dealing with matrices in REF form.
\end{remark}

\begin{example}
Use Gauss-Jordan Reduction to find the solution set of the following equations:
$$\begin{cases}
    2x + 3y + z = 2 \\
    2x + 4y + 2z = 8 \\
    x - y - z = 4
\end{cases}$$
\end{example}

\textit{Solution.} 

Form the augmented matrix and apply elementary row operations,
\begin{align*}
    \begin{amatrix}{3}
    2 & 3 & 1 & 2 \\ 
    2 & 4 & 2 & 8 \\
    1 & -1 & -1 & 4
\end{amatrix} &\xrightarrow{R_1 \leftrightarrow R_3}
\begin{amatrix}{3}
    1 & -1 & -1 & 4 \\
    2 & 4 & 2 & 8 \\
    2 & 3 & 1 & 2
\end{amatrix} \\
\xrightarrow{\substack{R_2 \leftarrow R_2 - R_1 \\ R_3 \leftarrow R_3 - R_1}}
\begin{amatrix}{3}
    1 & -1 & -1 & 4 \\ 
    0 & 6 & 4 & 0 \\
    0 & 5 & 3 & -6
\end{amatrix}
& \xrightarrow{R_2 \leftarrow R_2 - R_3}
\begin{amatrix}{3}
    1 & -1 & -1 & 4 \\
    0 & 1 & 1 & 6 \\
    0 & 5 & 3 & -6
\end{amatrix} \\
\xrightarrow{R_3 \leftarrow R_3 - 5R_2}
\begin{amatrix}{3}
    1 & -1 & -1 & 4 \\
    0 & 1 & 1 & 6 \\
    0 & 0 & -2 & -36
\end{amatrix} 
& \xrightarrow{R_3 \leftarrow -\frac{1}{2}R_3}
\begin{amatrix}{3}
    1 & -1 & -1 & 4 \\
    0 & 1 & 1 & 6 \\
    0 & 0 & 1 & 18
\end{amatrix} \\
\xrightarrow{R_1 \leftarrow R_1 + R_2}
\begin{amatrix}{3}
    1 & 0 & 0 & 10 \\
    0 & 1 & 1 & 6 \\
    0 & 0 & 1 & 18
\end{amatrix} 
& \xrightarrow{R_2 \leftarrow R_2 - R_3}
\begin{amatrix}{3}
    1 & 0 & 0 & 10 \\
    0 & 1 & 0 & -12 \\
    0 & 0 & 1 & 18
\end{amatrix} \\
\Longrightarrow &\begin{cases}
    x = 10 \\ 
    y = -12 \\
    z = 18
\end{cases}
\end{align*}

Therefore, the solution set is $\boxed{\{(10, -12, 18)\}}$.

\newpage
\begin{example}
Use Gauss-Jordan Reduction to find the solution set of the following equations:
$$\begin{cases}
    w + x + 2y + 3z = 13 \\
    w - 2x + y + z = 8 \\
    3w + x + y - z = 1
\end{cases}$$
\end{example}

\textit{Solution.} 

Form the augmented matrix and apply elementary row operations,
\begin{align*}
\begin{amatrix}{4}
    1 & 1 & 2 & 3 & 13 \\
    1 & -2 & 1 & 1 & 8 \\
    3 & 1 & 1 & -1 & 1
\end{amatrix} &\xrightarrow{\substack{R_2 \leftarrow R_2 - R_1 \\ R_3 \leftarrow R_3 - 3R_1}} 
\begin{amatrix}{4}
    1 & 1 & 2 & 3 & 13 \\
    0 & -3 & -1 & -2 & -5 \\
    0 & -2 & -5 & -10 & -38 
\end{amatrix} \\
\xrightarrow{R_2 \leftarrow R_2 - R_3} 
\begin{amatrix}{4}
    1 & 1 & 2 & 3 & 13 \\
    0 & -1 & 4 & 8 & 33 \\0 & -2 & -5 & -10 & -38 
\end{amatrix} & \xrightarrow{R_3 \leftarrow R_3 - 2R_2}
\begin{amatrix}{4}
    1 & 1 & 2 & 3 & 13 \\
    0 & -1 & 4 & 8 & 33 \\
    0 & 0 & -13 & -26 & -104
\end{amatrix} \\
\xrightarrow{\substack{R_2 \leftarrow -R_2 \\ R_3 \leftarrow -\frac{1}{13}R_3}} 
\begin{amatrix}{4}
    1 & 1 & 2 & 3 & 13 \\
    0 & 1 & -4 & -8 & -33 \\
    0 & 0 & 1 & 2 & 8
\end{amatrix} & \xrightarrow{R_1 \leftarrow R_1 - R_2} 
\begin{amatrix}{4}
    1 & 0 & 6 & 11 & 46 \\
    0 & 1 & -4 & -8 & -33 \\
    0 & 0 & 1 & 2 & 8
\end{amatrix} \\
\xrightarrow{\substack{R_1 \leftarrow R_1 - 6R_3 \\ R_2 \leftarrow R_2 + 4R_3}}
\begin{amatrix}{4}
    1 & 0 & 0 & -1 & -2 \\
    0 & 1 & 0 & 0 & -1 \\
    0 & 0 & 1 & 2 & 8
\end{amatrix} & \Longrightarrow 
\begin{cases}
    w - z = -2 \\
    x = -1 \\
    y + 2z = 8
\end{cases}
\end{align*}
Let $z = t$, then, $w = t - 2, y = 8-2t$. Therefore, the solution set is $\boxed{\{(t - 2, -1, -2t + 8, t) \mid t \in \RR\}}$.

\begin{remark}
\leavevmode
\begin{enumerate}
    \item The variable $z$ in the system of equations above is called a \textbf{free variable}.
    \item A system $Ax = b$ having at least one solution has $k$ free variables, where $k$ is the number of columns in the RREF of $[A \mid b]$ having no leading entries. The system is said to have a $k$-parameter solution.
\end{enumerate}
\end{remark}

For the rest of the examples, I will not show the whole process.\footnote{Since this is TOO tedious to format.} It is up to the reader to verify that the solution is correct.

\begin{example}
Use Gaussian Elimination to find the solution set.
$$\begin{cases}
    x - y + 2z = 4 \\
    x + z = 6 \\
    2x - 3y + 5z = 4 \\
    3x + 2y - z = 1
\end{cases}$$

\textit{Solution.}

Form the augmented matrix and apply elementary row operations,
\begin{align*}
    \begin{amatrix}{3}
        1 & -1 & 2 & 4 \\
        1 & 0 & 1 & 6 \\
        2 & -3 & 5 & 4 \\
        3 & 2 & -1 & 1
    \end{amatrix} &\xrightarrow{\text{REF}}
    \begin{amatrix}{3}
        1 & -1 & 2 & 4 \\
        0 & 1 & -1 & 2 \\
        0 & 0 & 1 & 21/2 \\
        0 & 0 & 0 & 1
    \end{amatrix} \\
\end{align*}
$$\Longrightarrow \begin{cases}
        x - y + 2z = 4 \\
        y - z = 2 \\
        z = 21/2 \\
        0 = 1 \Longrightarrow \text{This is a contradiction!}
    \end{cases}$$
Since the solution is inconsistent, therefore, the solution set is $\emptyset$ or $\{\}$.
\end{example}

\begin{remark}
The system $Ax = b$ has no solution iff $\bmat{0 & 0 & \cdots & 0 & \mid & 1}$ is a row in an REF of $[A \mid b]$.
\end{remark}

\begin{example}
Solve for the solution of the following system of equations.
$$\begin{cases}
    2x_1 + x_2 - 4x_3 - 5x_4 + 5x_5 = 5 \\
    3x_1 - x_2 - 6x_3 - 5x_4 = 5 \\
    5x_1 - 2x_2 - 10x_3 - 8x_4 - x_5 = 8
\end{cases}$$

\textit{Solution.} 

Form the augmented matrix and apply elementary row operations,
$$\begin{amatrix}{5}
    2 & 1 & -4 & -5 & 5 & 5 \\
    3 & -1 & -6 & -5 & 0 & 5 \\
    5 & -2 & -10 & -8 & -1 & 8
\end{amatrix} \xrightarrow{\text{RREF}}
\begin{amatrix}{5}
    1 & -2 & -2 & 0 & -5 & 0 \\
    0 & 1 & 0 & -1 & 3 & 1 \\
    0 & 0 & 0 & 0 & 0 & 0\\
\end{amatrix}
$$
$$\Longrightarrow \begin{cases}
    x_1 - 2x_2 - 2x_3 - 5x_5 = 0 \\
    x_2 - x_4 + 3x_5 = 1
\end{cases}$$
Let $x_3 = r, x_4 = s, x_5 = t$, then, $x_2 = s - 3t + 1$ and $x_1 = 2r + 2s - t + 2$. Therefore, the solution set is $\boxed{\{(2r + 2s - t + 2, s - 3t + 1, r, s,t) \mid r, s, t \in \RR\}}$.
\end{example}

\newpage
\subsection{Exercises}

\begin{problem}
Use the Gauss-Jordan reduction or Gaussian elimination to solve the system of equations
$$\begin{cases}
    3x + 4y - z  = 8 \\
    3x + 8y - 2z = 8 \\
    6x + 8y - 2z = 3
\end{cases}$$
\end{problem}

\begin{problem}
Let $a \in \mathbb{R}$, and consider the linear system in the unknowns $x_1, x_2, x_3, x_4$:
$$\begin{cases}
    x_1 +2x_2 -x_3 = 1 \\ 
    2x_1 +4x_2 -x_3 + x_4 = 1 \\
    x_1 +2x_2 +x_3 +2x_4 = \alpha
\end{cases}$$

(a) Find the value of $\alpha$ that would make the system consistent.

% \textit{Answer.} $\alpha = -1$

(b) When the system is consistent, find its solution set.

% \textit{Answer.} Solution Set: $\{(-2s - t, s, -t - 1, t) \mid s, t \in \mathbb{R}\}$
\end{problem}

\begin{problem}
Let
$$A = \bmat{\frac{1}{4}\sqrt{3} + \frac{1}{2} & \frac{1}{4}\sqrt{3} - \frac{1}{2} & -\frac{1}{4}\sqrt{2} \\ \frac{1}{4}\sqrt{3} - \frac{1}{2} & \frac{1}{4}\sqrt{3} - \frac{1}{2} & -\frac{1}{4}\sqrt{2} \\ \frac{1}{4}\sqrt{2} & \frac{1}{4}\sqrt{2} & \frac{1}{2}\sqrt{3}}$$
Find all $3 \times 1$ matrices $x$ satisfying $Ax = x$. [\textit{Hint: The matrix equation is equivalent to $(A - I_3)x = 0_{3 \times 1}$.}]
\end{problem}

\begin{problem}
Solve for the following system of nonlinear equations for $x, y$, and $z$.
\begin{align*}
x^2 + y^2 + z^2 = 6 \\
x^2 - y^2 + 2z^2 = 2 \\
2x^2 + y^2 - z^2 = 3
\end{align*}
[\textit{Hint: Let $X = x^2, Y = y^2, Z = z^2$.}]
\end{problem}

\begin{problem}
Solve the following system of nonlinear equations for the unknown angles $\alpha, \beta, \gamma$, where 
$0 \leq \alpha \leq 2\pi, 0 \leq \beta \leq 2\pi$, and $0 \leq \gamma \leq \pi$.
\begin{align*}
2\sin\alpha - \cos\beta + 3\tan\gamma = 3 \\
4\sin\alpha + 2\cos\beta - 2\tan\gamma = 2 \\
6\sin\alpha - 3\cos\beta + \tan\gamma = 9
\end{align*}
\end{problem}
\newpage
\section{Inverse of a Square Matrix}
Recall that for all $a \in \mathbb{R}$, there exists a multiplicative inverse $a^{-1} \in \mathbb{R}$, such that 
$$a \cdot a^{-1} = 1 =  a^{-1} \cdot a$$

Now let us figure out if $0_{n \times n} + A \in M_n(\mathbb{R})$, can we find a $B \in M_n(\mathbb{R})$ such that $AB = I_n = BA$? The answer to that is \textbf{not always}!

As a counterexample, suppose that we have two matrices $A$ and $B$, such that their product $AB = I_2$,
$$A = \left [
\begin{matrix}
1 & 0 \\
0 & 0 \\
\end{matrix} \right ] \hspace{0.3cm} \text{and} \hspace{0.3cm} 
B = \left [
\begin{matrix}
a & b \\
c & d \\
\end{matrix} \right ]$$

Now suppose this is true, then,
$$\left [
\begin{matrix}
1 & 0 \\
0 & 0 \\
\end{matrix} \right ]
\left [
\begin{matrix}
a & b \\
c & d \\
\end{matrix} \right ] = 
\left [
\begin{matrix}
1 & 0 \\
0 & 1 \\
\end{matrix} \right ]$$
$$\left [
\begin{matrix}
a & b \\
0 & 0 \\
\end{matrix} \right ] = 
\left [
\begin{matrix}
1 & 0 \\
0 & 1 \\
\end{matrix} \right ]$$

This is a contradiction, because if we compare the $(2, 2)$ entries of both matrices, $0 \neq 1$. Hence no such $B$ exists.

\begin{definition}
A square matrix $A \in M_n(\mathbb{R})$ is said to be \textbf{nonsingular/invertible} if there exists a matrix $B \in M_n(\mathbb{R})$ such that
$$AB = I_n = BA$$
In this case, $B$ is said to be the multiplicative inverse of $A$. If $A$ has no inverse, it is said to be \textbf{singular/invertible}.
\end{definition}

\begin{theorem}
The inverse of a nonsingular matrix is unique.
\end{theorem}

\begin{proof}
Assume for the sake of contradiction that $A$ has two inverse $B$ and $C$, such that
$$AB = BA = I_n$$
$$AC = CA = I_n$$
Recall that $B$ can be expressed as
$$B = BI_n = B(AC)$$
$$= (BA)C = I_nC = C$$
Because $B = C$, which is a contradiction, therefore, the inverse is unique.

\end{proof}

\textbf{Notation:} Let $A \in M_n(\mathbb{R})$,

1. If $A$ is nonsingular, then its unique inverse is denoted by $A^{-1}$

2. If $A \neq 0_n$ then $A^0 = I_n$

3. For any $k \in \mathbb{Z}^+$, we represented the matrix being multiplied $k$ times as
$$A^k = \underbrace{A \times A \times \cdots \times A}_{k}$$

4. The laws of exponents also apply,
$$A^k \cdot A^l = A^{k+l}$$
$$(A^k)^l = A^{kl}$$
However, 
$$(AB)^2 \neq A^2B^2$$

\begin{remark}(Properties of Nonsingular Matrices).

Let $A, B \in M_n(\mathbb{R})$ and nonsingular, $0 \in \mathbb{R}$, and $k \in \mathbb{Z}^+$. Then the following are nonsingular:
\begin{enumerate}
    \item $A$ with its inverse $A^{-1}$
    \item $AB$ with its inverse\footnote{This is also known as the \href{https://math.stackexchange.com/questions/2435303/russells-shoes-and-socks}{socks-shoes principle}.} $(AB)^{-1} = B^{-1}A^{-1}$
    \item $cA$ with its inverse $\displaystyle \frac{1}{c} \cdot A^{-1}$
    \item $A^{\mathsf{T}}$ with its inverse $(A^{\mathsf{T}})^{-1} = (A^{-1})^{\mathsf{T}}$
    \item $A^k$ with its inverse $(A^k)^{-1} = (A^{-1})^k$
\end{enumerate}
\end{remark}

\begin{theorem}[Cancellation Laws]
Let $B$ and $C$ be matrices of suitable sizes and $A \in M_n(\mathbb{R})$ be nonsingular. Then,
$$AB = AC \Longrightarrow B = C$$
$$BA = CA \Longrightarrow B = C$$
\end{theorem}

\begin{proof}
Since $A$ is nonsingular and its inverse $A^{-1}$ exists, then suppose that $AB = AC$ is true, we can show that by multiplying both sides by its inverse,
$$A^{-1}(AB) = A^{-1}(AC)$$
$$I_nB = I_nC$$
$$B = C$$

Similarly we can show that $B = C$ when $BA = CA$,
$$(BA)A^{-1} = (CA)A^{-1}$$
$$BI_n = CI_n$$
$$B = C$$

\end{proof}

\begin{remark}
If $A, B \in M_n(\mathbb{R})$ such that $AB = I_n$, then $BA = I_n$ and so $A$ is nonsingular.
\end{remark}

To further understand this remark, let us put it into an example. Suppose $A \in M_n(\mathbb{R})$ such that $A^2 - 2A + I_2 = 0$, figure out whether $A$ is nonsingular.

\textit{Solution.} 
$$I_2 = 2A - A^2$$
$$=(2I_2)(A)- AA$$
$$=\underbrace{(2I_2 - A)}_B(A)$$

Therefore, $A$ is nonsingular with inverse $A^{-1} = 2I_2 - A$.

\newpage
\begin{fact}
If $A$ and $B$ are matrices of  size and $b_j = \text{col }j$ of $B$, then
$$AB = A\left[\begin{matrix}
    b_1 & b_2 & \cdots & b_n
\end{matrix}\right] = \left[
\begin{matrix}
    Ab_1 & Ab_2 & \cdots & Ab_n
\end{matrix}
\right]$$
\end{fact}

\begin{theorem}
$A \in M_n(\mathbb{R})$ is nonsingular iff the system of equations, $Ax = b$ has a unique solution for any $b$, given by 
$$x = A^{-1}b$$
\end{theorem}

\begin{proof}
Suppose that $A$ is nonsingular and let $b \in M_{n \times 1}(\mathbb{R})$. Take $x = A^{-1}b$,
$$A(A^{-1}b) = I_nb = b$$
Therefore, $Ax = b$ has at least one solution. Let us use proof by contradiction to show that the system has a unique solution. Suppose then that $x_1$ and $x_2$ are solutions of $Ax = b$ such that
$$Ax_1 = b = Ax_2$$
By cancellation law, we can show that, $x_1 = x_2$. Which is a contradiction, hence, the system $Ax = b$ has only one solution. Suppose $Ax = b$ has exactly one solution for any $b \in M_{n \times 1}(\mathbb{R})$, if $e_j = \text{col}\, j$ of $I_n$, then $Ax = e_j$ has a unique solution $x_j$.

Let $B = \bmat{x_1 & x_2 & \cdots & x_n}$,
\begin{align*}
AB = A\bmat{x_1 & x_2 \cdots & x_n} \\
= \bmat{Ax_1 & Ax_2 & \cdots & Ax_n} \\
= \bmat{e_1 & e_2 \cdots & e_n} = I_n
\end{align*}

Therefore, $A$ is nonsingular with inverse $B$.
\end{proof}

\begin{remark}
Let $A \in M_n(\mathbb{R})$ be nonsingular. Then column $j$ of $A^{-1}$ is the unique solution to $Ax = e_j$ where $e_j$ is column $j$ of $I_n$.
\end{remark}

\lline

We will abbreviate the elementary row operations in some texts as \textbf{EROs}.

\begin{itemize}
    \item $A \in M_n(\mathbb{R})$ is nonsingular if there exists an inverse $A^{-1} \in M_n(\mathbb{R})$ such that
    $$A\cdot A^{-1} = I_n = A^{-1} \cdot A$$
    \item To find $A^{-1}$
    $$[A \mid I_n] \stackrel{\text{EROs}}{\longrightarrow} [C \mid D]$$
    wherein $C = I_n$ and $D = A^{-1}$
    \item If $C$ has a row of 0s, then $A$ is singular

\end{itemize}

\begin{remark}
$A \in M_n(\mathbb{R})$ is nonsingular iff $A$ is row equivalent to $I_n$, and $A$ is said to be singular iff $A$ is row equivalent to a matrix of zero row.
\end{remark}

\begin{definition}
A matrix $E \in M_n(\mathbb{R})$ is called an \textbf{elementary matrix} if it can be obtained by applying a single elementary row operation on $I_n$.
\end{definition}

\newpage
\begin{example}\label{example:3.1}
Let $A = \bmat{2 & -4 & 6 \\ 3 & -5 & 8 \\ 0 & 1 & -4}$, use elementary row operations to find the inverse $A\inv$.
\end{example}

\textit{Solution.} 

Form the augmented matrix,
$$\bmat{2 & -4 & 6 & \vline & 1 & 0 & 0 \\
3 & -5 & 8 & \vline & 0 & 1 & 0 \\ 0 & 1 & -4 & \vline & 0 & 0 & 1} \xrightarrow{\text{RREF}} 
\bmat{1 & 0 & 0 & \vline & -2 & 5/3 & 1/3 \\ 0 & 1 & 0 & \vline & -2 & 4/3 & -1/3 \\ 0 & 0 & 1 & \vline & -1/2 & 1/3 & -1/3}
$$
$$\therefore A\inv = \bmat{-2 & 5/3 & 1/3 \\ -2 & 4/3 & -1/3 \\ -1/2 & 1/3 & -1/3}$$

\begin{example}
Use the answer from \hyperref[example:3.1]{Example 3.1} to find the unique solution of the system of equations:
$$\begin{cases}
    2x - 4y + 6z = 1 \\
    3x - 5y + 8z = -2 \\
    y - 4z = 4
\end{cases}$$

\textit{Solution.}

Recall that $x = A\inv b$,
$$\bmat{x \\ y \\ z} = \bmat{-2 & 5/3 & 1/3 \\ -2 & 4/3 & -1/3 \\ -1/2 & 1/3 & -1/3}\bmat{1 \\ -2 \\ 4}$$
$$\therefore \bmat{x \\ y \\ z } = \bmat{-4 \\ -6 \\ -5/2}$$
\end{example}

\begin{remark}
$A \in M_n(\RR)$ is nonsingular iff $A$ is row equivalent to $I_n$. And $A$ is singular iff $A$ is row equivalent to a matrix of zero row.
\end{remark}

\begin{definition}
A matrix $E \in M_n(\RR)$ is called an elementary matrix if it can be obtained by applying a single elementary row operation on $I_n$.
\end{definition}

\underline{\textbf{Types of Elementary Matrices}}

\underline{\textbf{Type 1}}. We denote the elementary matrix obtained when we switch rows $k$ and $\ell$ of $I_n$ as $P_{k\ell}$

e.g. $\displaystyle \bmat{0 & 0 & 0 & 1 \\ 0 & 1 & 0 & 0 \\ 0 & 0 & 1 & 0 \\ 1 & 0 & 0 & 0} = P_{14}; \hspace{0.2cm} P_{14}^{-1} = P_{14}$

\underline{\textbf{Type 2}}. We denote the elementary matrix obtained when we multiply row $k$ of $I_n$ by the nonzero constant $c$ as $Q_k(c)$

e.g. $\displaystyle \bmat{1 & 0 & 0 & 0 \\ 0 & 1 & 0 & 0 \\ 0 & 0 & c & 0 \\ 0 & 0 & 0 & 1} = Q_3(c); \hspace{0.2cm} Q_3(c)^{-1} = Q_3\left(\frac{1}{c}\right)$

\underline{\textbf{Type 3}}. We denote the elementary matrix obtained when we add $d$ times row $k$ of $I_n$ to row $\ell$

e.g. $\displaystyle \bmat{1 & d & 0 & 0 \\ 0 & 1 & 0 & 0 \\ 0 & 0 & 1 & 0 \\ 0 & 0 & 0 & 1} = R_{21}(d); \hspace{0.2cm} R_{21}(d)^{-1} = R_{21}(-d)$

\begin{theorem}
The elementary matrices are nonsingular and their inverses are also elementary.
\end{theorem}

\begin{remark}
Some notable observations for elementary matrices are shown here:

1. The inverses of the elementary matrices can be simplified as follows:
\begin{itemize}
    \item $P_{k\ell}^{-1} = P_{k\ell}$
    \item $\displaystyle Q_k(c)^{-1} = Q_k\left(\frac{1}{c}\right), c \neq 0$
    \item $R_{k\ell}(c)^{-1} = R_{k\ell}(-c)$
\end{itemize}

2. Let $E \in M_n(\mathbb{R})$ be an elementary matrix and $C \in M_{n \times p}(\mathbb{R})$. Then, we can denote the matrix obtained by applying $C$ to the same row operation applied on $I_n$ to get $E$ as $EC$

\end{remark}

\begin{example}\label{example:3.3}
Let $E_1, E_2, E_3$, and $A$ be matrices defined as follows:
\begin{align*}
E_1 &= \begin{bmatrix} 1 & 0 & 0 \\ 0 & \sqrt{2} & 0 \\ 0 & 0 & 1 \end{bmatrix}, &
E_2 &= \begin{bmatrix} 1 & 0 & 0 \\ 0 & 1 & 4 \\ 0 & 0 & 1 \end{bmatrix}, \\[1ex]
E_3 &= \begin{bmatrix} 1 & 0 & 0 \\ 0 & 0 & 1 \\ 0 & 1 & 0 \end{bmatrix}, &
A   &= \begin{bmatrix} 1 & \pi & 3 & 2 \\ 0 & \sqrt{2} & 0 & 0 \\ 0 & -\frac{3}{4} & -1 & 5 \end{bmatrix}
\end{align*}
Determine the elementary row operations corresponding to $E_1$, $E_2$, and $E_3$.

\textit{Solution.} The elementary row operations represented by the given matrices are:
\[
\begin{aligned}
E_1 &\implies R_2 \leftarrow \sqrt{2}R_2 && \text{(or $Q_2(\sqrt{2})$)} \\
E_2 &\implies R_2 \leftarrow R_2 + 4R_3 && \text{(or $R_{32}(4)$)} \\
E_3 &\implies R_2 \longleftrightarrow R_3 && \text{(or $P_{23}$)}
\end{aligned}
\]
\end{example}

\newpage
\begin{example}
Find $E_3E_2E_1A$ using the matrices from \hyperref[example:3.3]{Example 3.3}.
\end{example}

\textit{Solution.} 
\[
\begin{aligned}
E_3E_2E_1A = E_3E_2E_1\bmat{1& \pi & 3 & 2 \\ 0 & -\sqrt{2} & 0 & 0 \\ 0 & -3/4 & -1 & 5} \\
= E_3E_2\bmat{1 & \pi & 3 & 2 \\ 0 & -2 & 0 & 0 \\ 0 & -3/4 & -1 & 5} \\
=E_3\bmat{1 & \pi & 3 & 2 \\ 0 & -5 & -4 & 20 \\ 0 & -3/4 & -1 & 5} \\
=\bmat{1  & \pi & 3 & 2 \\ 0 & -3/4 & -1 & 5 \\ 0 & -5 & -4 & 20}
\end{aligned}
\]

\begin{theorem}
Let $C, D \in M_{m \times n}(\mathbb{R})$. Then $C$ and $D$ are row equivalent iff there exists elementary matrices $E_1, E_2, \dots, E_n \in M_n(\mathbb{R})$ such that
$$D = E_k \cdots E_2E_1C$$
\end{theorem}

\begin{remark}
$A \in M_n(\mathbb{R})$ is nonsingular iff:
\begin{itemize}
    \item $A$ is row equivalent to $I_n$
    \item $I_n = E_k \cdots E_2E_1A$ for some elementary $E_k, \dots, E_2, E_1 \in M_n(\mathbb{R})$
    \item $A^{-1} = E_k \cdots E_2E_1$ and $A = E_1^{-1}E_2^{-1}\cdots E_k^{-1}$ for some elementary $E_1, E_2, \dots, E_k \in M_n(\mathbb{R})$
    \item $A$ is a product of elementary matrices
\end{itemize}
\end{remark}

\begin{example}
Let $A = \bmat{ 3 & -2 \\ 1 & 0}$. Write $A$ and $A^{-1}$ as products of elementary matrices.
\end{example}

\textit{Solution.} Let us first reduce $A$ into the identity matrix and then identify the EROs used,
\begin{align*}
\bmat{3 & -2 \\ 1 & 0} \xrightarrow{R_1 \leftrightarrow R_2}
\bmat{1 & 0 \\ 3 & -2} \xrightarrow{R_2 \to R_3 - 3R_2}
\bmat{1 & 0 \\ 0 & -2} \\
\bmat{1 & 0 \\ 0 & -2} \xrightarrow{R_2 \to -1/2R_2} \bmat{1 & 0 \\ 0 & 1} = I_2
\end{align*}
Then,
$$I_2 = Q_2\left(-\frac{1}{2}\right)R_{12}(-3)P_{12}A$$
Its inverse $A^{-1}$ is:
\begin{align*}
A^{-1} &= Q_2\left(-\frac{1}{2}\right)R_{12}(-3)P_{12} \\
&=\bmat{1 & 0 \\ 0 & -1/2}\bmat{1 & 0 \\ -3 & 1}\bmat{0 & 1 \\ 1 & 0}
\end{align*}

Recall that $A = E_1^{-1}E_2^{-1}\cdots E_k^{-1}$, thus,
$$A = \boxed{\bmat{0 & 1 \\ 1 & 0}\bmat{1 & 0 \\ 3 & 1}\bmat{1 & 0 \\ 0 & -2}}$$

\newpage
\begin{theorem}
Let $A \in M_n(\mathbb{R})$. The following statements are equivalent:
\begin{itemize}
    \item $A$ is nonsingular
    \item $A$ is row equivalent to $I_n$
    \item $A$ is the product of elementary matrices
    \item The matrix equation $Ax = b$ has a unique solution for any $b \in M_{n \times 1}(\mathbb{R})$
    \item The homogeneous matrix equation $Ax = 0_{n \times 1}$, has only the \underline{\textbf{trivial solution}} $x = 0_{n \times 1}$
\end{itemize}
\end{theorem}

\begin{example}
$A = \bmat{1 & 1 & 1 \\ 1 & 1 & 1 \\ 1 & 1 & 1}$ is singular, since $Ax = 0_{3 \times 1}$ has a nontrivial solution namely $x = \bmat{1 \\ 1 \\ -2}$.
\end{example}
\newpage
\section{Determinant of a Square Matrix}
Let $A = [a_{ij}]\in M_n(\mathbb{R})$. The determinant of $A$ is defined as follows:
\begin{itemize}
    \item If $n = 1$, then $\det A = a_{11}$
    \item Suppose $n \geq 1$ and assume that the determinant of a matrix of size $n - 1$ is defined. Let $M_{ij}$ be the matrix obtained from $A$ by deleting row $i$ and column $j$ of $A$. Define
    $$C_{ij} = (-1)^{i+j} \det M_{ij}$$ wherein $(i, j)$ is the cofactor of $A$.
\end{itemize}

Then we define the determinant as a cofactor expansion along row $i$,
$$\det A = a_{i1}C_{i1} + a_{i2}C_{i2} + \cdots + a_{in}C_{in}$$

Similarly, we could also define it as a cofactor expansion along column $j$,
$$\det A = a_{1j}C_{1j} + a_{2j}C_{2j} + \cdots + a_{nj}C_{nj}$$

\lline

\begin{example}
Obtain a formula for the determinant of $\displaystyle A = \left[ 
\begin{matrix}
    a_{11} & a_{12} \\
    a_{21} & a_{22}
\end{matrix}
\right]$
\end{example}

\textit{Solution.} By cofactor expansion along row 1,
\begin{align*}
\det A &= a_{11}C_{11} + a_{12}C_{12} \\
&=a_{11}(-1)^{1+1}\det [a_{22}] + a_{12}(-1)^{1+2}\det [a_{21}] \\
&=\boxed{a_{11}a_{22} - a_{12}a_{21}}
\end{align*}

\begin{example}
Find the determinant of $\displaystyle B = \left[
\begin{matrix}
    5 & 6 \\
    4 & 8
\end{matrix}
\right]$.
\end{example}

\textit{Solution.} Recall that from the previous example, we have obtained a determinant formula for a square matrix of size 2, we could use that to solve for $\det B$
$$\det B = 5(8) - 6(4) = \boxed{16}$$

\begin{example}
Find the determinant of $\displaystyle B = \left[ 
\begin{matrix}
    3 & 1 & -4 \\
    2 & 5 & 6 \\
    1 & 4 & 8
\end{matrix}
\right]$.
\end{example}

\textit{Solution.} By cofactor expansion along column 1,
$$\det B = 3(-1)^{1+1}
\vmat{
5 & 6 \\
4 & 8
} + 
2(-1)^{2 + 1}
\vmat{
1 & -4 \\
4 & 8
} +
1(-1)^{3 + 1}
\vmat{
1 & -4 \\
5 & 6
}
$$
$$=3(16) -2(24) + 26 = \boxed{26}$$

\newpage
\begin{example}
Find the determinant of $\displaystyle C = \bmat{
2 & 3 & 0 & 0 \\
1 & 1 & 0 & 0 \\
4 & 6 & 1 & -5 \\
0 & 5 & 3 & 0
}$.
\end{example}

\textit{Solution.} Notice that column 4 has 3 zeros which could help us compute for the determinant for efficiently,
$$\det C = 0 + 0 -5(-1)^{3 + 4}\vmat{
2 & 3 & 0 \\
1 & 1 & 0 \\
0 & 5 & 3
} + 0 = 5 \vmat{
2 & 3 & 0 \\
1 & 1 & 0 \\
0 & 5 & 3
}$$
We will perform cofactor expansion along column 3 since it has the least amount of zeros,
$$=5\left(3(-1)^{3+3}\vmat{
2 & 3 \\
1 & 1
}\right) = 15(2 - 3) = \boxed{-15}$$

\begin{example}
The examples that will be shown here will help us determine some properties of determinants. Let $\displaystyle A = \bmat{
a & b \\
c & d
}$, and $\displaystyle B = \bmat{
a_{11} & a_{12} & a_{13} \\
0 & a_{22} & a_{23} \\
0 & 0 & a_{33}
}$
\begin{itemize}
    \item Suppose we switch row 1 and row 2 of matrix $A$, then the determinant is
    $$\vmat{
    c & d \\
    a & b
    } = bc - ad = -(ad -bc) = -\det A$$
    \item Suppose we add a constant $r \in \mathbb{R}$ to row 2 of matrix $A$, then the determinant is
    $$\vmat{
    a & b\\
    rc & rd
    } = ard - brc = r(ad - bc) = r\det A$$
    \item Suppose we multiply $r$ to row 1 of matrix $A$ and then add it to row 2, then the determinant is
    $$\vmat{
    a & b\\
    ra + c & rb + d
    } = a(rb + d) - b(ra + c) = ad - bc = \det A$$
    \item Suppose we perform cofactor expansion along row 3 of matrix $B$ to get its determinant,
    $$\det B = a_{33}(-1)^{3 + 3}\vmat{
    a_{11} & a_{12} \\
    0 & a_{22}
    } = a_{33}(a_{11}a_{22}) = a_{11}a_{22}a_{33}$$
\end{itemize}
\end{example}

\lline

\underline{\textbf{Some Properties of Determinant}}. Let $A, B \in M_n(\mathbb{R})$,

1. If $B$ is obtained from $A$ by:
\begin{itemize}
    \item switching two rows/columns of $A$ then $\det B = -\det A$
    \item multiplying $r \in \mathbb{R}$ to a row/columns of $A$ then $\det B = r\det A$

    e.g. $\det(rA) = \underbrace{r \times r \times \cdots \times r}_{n}\det A = r^n\det A$
    \item adding a multiple of a row/column of $A$ to another row/column of $A$ then $\det B = \det A$
\end{itemize}

2. $\det(A^{\mathsf{T}}) = \det A$

3. If $A$ has a row/column full of 0s then $\det A = 0$.

4. If $A$ has two equal rows/columns, then $\det A = 0$.

5. If $A = [a_{ij}]$ is upper or lower triangular, then $\det A = a_{11}a_{22}\cdots a_{nn}$, which is the product of all the main diagonal entries of the matrix.

6. $\det(AB) = \det A \det B$, and in particular, $\det(A^k) = (\det A)^k$ for all $k \in \mathbb{Z}^+$.

\lline

\begin{theorem}
If $E \in M_n(\mathbb{R})$ is elementary, then $\det E \neq 0$.
\end{theorem}

\begin{theorem}
$A \in M_n(\mathbb{R})$ is nonsingular iff $\det A \neq 0$, in this case
$$\det(A^{-1}) = \frac{1}{\det A}$$
\end{theorem}

\begin{definition}
Let $A \in M_n(\mathbb{R})$ and $C_{ij}$ be the $(i, j)$ column of $A$. The \textbf{adjoint} of $A$ is $\adj A = [C_{ij}]^{\mathsf{T}}$.
\end{definition}

\begin{example}
Find the adjoint of the following:

1. $\displaystyle A = \bmat{a & b \\ c & d}$

\textit{Solution.}

\begin{align*}
C_{11} = (-1)^{1 + 1}d = d \quad & C_{21} = (-1)^{2 + 1} b = -b \\
C_{12} = (-1)^{1 + 2}c = -c \quad & C_{22} = (-1)^{2 + 2}a = a
\end{align*}

$$\therefore \adj A = \bmat{ d & -b \\ -c & a}$$

2. $\displaystyle B = \bmat{1 & 1 & 0 \\ 2 & 3 & 4 \\ 5 & 8 & 9}$
\end{example}

\textit{Solution.} \label{example:4.6}
\begin{align*}
C_{11} &= \vmat{3 & 4 \\ 8 & 9} = -5 & C_{21} &= -\vmat{1 & 0 \\ 8 & 9} = -9 & C_{31} &= \vmat{1 & 0 \\ 3 & 4} = 4 \\
C_{12} &= -\vmat{2 & 4 \\ 5 & 9} = 2 & C_{22} &= \vmat{1 & 0 \\ 5 & 9} = 9  & C_{32} &= -\vmat{1 & 0 \\ 2 & 4} = -4 \\
C_{13} &= \vmat{2 & 3 \\ 5 & 8} = 1  & C_{23} &= -\vmat{1 & 1 \\ 5 & 8} = -3 & C_{33} &= \vmat{1 & 1 \\ 2 & 3} = 1
\end{align*}
$$\therefore \adj B = \bmat{-5 & -9 & 4 \\ 2 & 9 & -4 \\ 1 & -3 & 1}$$

\begin{theorem}[Adjoint Formula]
If $A \in M_n(\mathbb{R})$ then,
$$A(\adj A) = (\det A)I_n = (\adj A)A$$
If $A$ is nonsingular then,
$$A\left(\frac{1}{\det A} \cdot \adj A\right) = I_n$$
$$A^{-1} = \frac{1}{\det A} \adj A$$
\end{theorem}

\begin{example}
If $\displaystyle B = \bmat{
a & b \\
c & d
}$ is nonsingular then,
$$A^{-1} = \frac{1}{ad - bc}\bmat{d & -b \\ -c & a}$$
\end{example}

\begin{example}
Find the inverse of $\displaystyle B = \bmat{1 & 1 & 0 \\ 2 & 3 & 4 \\ 5 & 8 & 9}$.
\end{example}

\textit{Solution.}

Recall that if we perform Type III operations in any row or column of a matrix, then its determinant does not change. Thus, we could multiply -1 to the the first column and add that to the second column, we would get the following matrix:
$$\bmat{1 & 0 & 0 \\ 2 &1 & 4 \\ 5 & 3 & 9} \Longrightarrow \det B = \vmat{1 & 0 & 0 \\ 2 & 1 & 4 \\ 5 & 3 & 9}$$

Let us perform cofactor expansion along row 1,
$$\det B = 1 \vmat{1 & 4 \\ 3 & 9} = 9 - 12 = -3$$

Since we have already obtained $\adj B$ from \hyperref[example:4.6]{Example 4.6}, we could use that to find the inverse using the adjoint formula,
$$B^{-1} = -\frac{1}{3}\bmat{-5 & -9 & 4 \\ 2 & 9 & -4 \\ 1 & -3 & 1}=\boxed{\bmat{5/3 & 3 & -4/3 \\ -2/3 & -3 & 4/3 \\ -1/3 & 1 & -1/3}}$$

\begin{theorem}[Cramer's Rule]
Let $A \in M_n(\mathbb{R})$ be nonsingular and $b \in M_{n \times 1}(\mathbb{R})$. The matrix equation $Ax = b$ has a unique solution $\displaystyle x = \bmat{x_1 & x_2 & \cdots & x_n}^{\mathsf{T}}$ given by
$$x_i = \frac{\det A_i}{\det A}$$
where $A_i$ is the matrix obtained by replacing column $i$ of $A$ by $b$.
\end{theorem}

\begin{example}
Solve for $y$ only in the system of equations below
$$\begin{cases}
    x + y + z = 8 \\
    -x - 2y + 3z = 1 \\
    3x - 7y + 4z = 0
\end{cases}$$
\end{example}

\textit{Solution.}

Let $A$ be the coefficient matrix and $b$ be the column vector of the constants,
$$A = \bmat{
1 & 1 & 1 \\
-1 & -2 & 3 \\
3 & -7 & 4
}, \hspace{0.3cm} b = \bmat{
8 \\ 1 \\ 0
}, \hspace{0.3cm} A_2 = \bmat{1 & 8 & 1 \\ -1 & 1 & 3 \\ 3 & 0 & 4}
$$
By Cramer's Rule,
$$y = \frac{\det A_2}{\det A} = \frac{105}{39} = \boxed{\frac{35}{13}}$$
\newpage
\section{Vector Spaces and Subspaces}

\begin{definition}
A \textbf{vector space} (over $\mathbb{R}$) is a set $V$ equipped with two operations: $\oplus$ (vector addition) and $\odot$ (scalar multiplication) satisfying the following for all $u, v, w \in V$ and for all $c, d \in \mathbb{R}$:
\begin{itemize}
    \item A1. $u \oplus v \in V$ (closure)
    \item A2. $(u \oplus v) \oplus w = u \oplus (v \oplus w)$ (associativity)
    \item A3. $u \oplus v = v \oplus u$ (commutativity)
    \item A4. There exists a zero vector $\vec{0} \in V$ such that $u \oplus \vec{0} = u$
    \item A4. For any $u \in V$, there exists an inverse $-u \in V$ such that $u \oplus (-u) = \vec{0}$ \\
    \leavevmode
    \item S1. $c \odot u \in V$ (closure)
    \item S2. $(c + d) \odot u = (c \odot u) \oplus (d \odot u)$ 
    \item S3. $c \odot (u \oplus v) = (c \odot u) \oplus (c \odot v)$
    \item S4. $(cd) \odot u = c \odot (d \odot u)$
    \item S5. $1 \odot u = u$
\end{itemize}
\end{definition}

\begin{remark}
Let $V$ be a vector space
\begin{enumerate}
    \item The elements of $V$ are called \textbf{vectors}, and the set of all real numbers are \textbf{scalars}.
    \item There are 2 different:
    \begin{enumerate}
        \item additions: $\oplus$ on $V$ and $+$ on $\RR$
        \item multiplications: $\odot$ on $V$ and $\cdot$ on $\RR$
        \item zeros: $\vec{0}$ (zero vector) and $0$ (zero scalar)
    \end{enumerate}
\end{enumerate}
If the context is clear we will use the $+$ and $\cdot$ in place of $\oplus$ or $\odot$.
\end{remark}

\begin{example}
The following are vector spaces:

1. $\mathbb{R}^n = \{(a_1, a_2, \dots, a_n) \mid a_1, a_2, \dots, a_n \in \mathbb{R})\}$ under 
$$(a_1, a_2, \dots, a_n) + (b_1, b_2, \dots, b_n) = (a_1 + b_1, a_2 + b_2, \dots, a_n + b_n)$$
$$c(a_1, a_2, \dots, a_n) = (ca_1, ca_2, \dots, ca_n)$$

Wherein the zero vector $\vec{0} = (\underbrace{0, 0, \dots, 0}_n)$ and its inverse is $-(a_1, a_2, \dots, a_n) = (-a_1, -a_2, \dots, -a_n)$

2. $M_{m \times n}(\mathbb{R})$ under matrix addition and scalar multiplication on a matrix

3. $P_n = \{a_nx^n + a_{n-1}x^{n-1} + \dots + a_1x + a_0 \mid a_0, a_1, \dots, a_n \in \mathbb{R}\}$

4. $\mathbb{C} = \{a + bi \mid a, b \in \mathbb{R}\, \land \, i^2 = -1\}$

5. $\mathscr{F}(\mathbb{R})$ which is the set of all real-valued functions with domain $\mathbb{R}$ under
$$(f + g)(x) = f(x) + g(x)$$
$$(cf)(x) = cf(x) , \quad c \in \RR$$

6. $V = (0, +\infty) = \mathbb{R}^+$ under
$$x \oplus y = xy \in V$$ 
$$r \odot x = x^r \in V$$


\end{example}

\begin{theorem}
Let $V$ be a vector space, $u \in V$ and $c \in \mathbb{R}$.Then
\begin{itemize}
    \item $\vec{0}$ is unique and the additive inverse $-u$ is unique.
    \item $c\vec{0} = c(\vec{0} + \vec{0})$
    \item $cu = \vec{0} \Longrightarrow c = 0 \text{ or } u = \vec{0}$
    \item $-(cu) = (-c)u = c(-u)$
    \item $-u = (-1)u$
\end{itemize}
\end{theorem}

\begin{definition}
Let $V$ be a vector space and $W \subseteq V$. If $W$ is also a vector space under the same operations as $V$, then $W$ is called a \textbf{subspace} of $V$, denoted by $W \leq V$.
\end{definition}

\begin{example}
\leavevmode
\begin{enumerate}
    \item If $V$ is a vector space then the following are subspaces of $V$:
    \begin{enumerate}
        \item $V$ (improper subspace)
        \item $\{\vec{0}\}$ (trivial subspace)
    \end{enumerate}
    \item $P_1 = \{ax + b \mid a, b \in \RR\}$ is a subspace of $P_2 = \{ax^2 + bx + c \mid a, b, c \in \RR\}$.
\end{enumerate}
\end{example}

\begin{theorem}[Subspace Test]
Let $V$ be a vector space and $W \subseteq V$. Then $W \leq V$ iff the following are satisfied:
\begin{itemize}
    \item $W \neq \emptyset$
    \item For all $w_1, w_2 \in W$, $w_1 + w_2 \in W$
    \item For all $c \in \mathbb{R}$, and for all $w \in W$, $cw \in W$
\end{itemize}
\end{theorem}

\begin{remark}
A nonempty subset $W$ of $V$ is \ul{NOT} a subspace of $V$ iff at least one of the following is true:
\begin{enumerate}
    \item There exists $w_1, w_2 \in W$ such that $w_1 + w_2 \notin W$
    \item There exists $c \in \RR$ such that $cw \notin W$
\end{enumerate}
\end{remark}

\begin{example}
Let $V = \RR^3$. Show that $W = \{(x, y, z) \mid x + y + z = 0\}$ is a subspace of $\RR^3$.
\end{example}

\begin{proof}
\leavevmode
\begin{itemize}
    \item $(0, 0, 0) \in W$ since $0 + 0 + 0 = 0$, thus, $W \neq \emptyset$.
    \item Let $w_1 = (x_1, y_1, z_1), w_2 = (x_2, y_2, z_2) \in W$
    $$w_1 + w_2 = (x_1 + x_2, y_1 + y_2, z_1 + z_2) \in W$$
    \begin{align*}
        (x_1 + x_2) + (y_1 + y_2) + (z_1 + z_2) &= (x_1 + y_1 + z_1) + (x_2 + y_2 + z_2) \\
        &= 0
    \end{align*}
    Thus, $w_1 + w_2 \in W$
    \item Let $c \in \RR$ and $w = (x, y, z) \in W$
    $$c(x, y, z) = (cx, cy, cz) $$
    $$c(x + y + z) = c(0)$$
    Thus, $cw \in W$
\end{itemize}
Therefore, $W \leq \RR^3$ by the Subspace Test.
\end{proof}

\begin{example}
Let $V = M_2(\RR)$. Show that 
$$W = \{A \in M_2(\RR) \mid A \text{ is upper triangular}\}$$
is a subspace of $V$.
\end{example}

\begin{proof}
\leavevmode
\begin{itemize}
    \item $0_{2 \times 2} \in W \Longrightarrow W \neq \emptyset$
    \item Let $A = \bmat{a_1 & a_2 \\ 0 & a_3}$, $B = \bmat{b_1 & b_2 \\ 0 & b_3}$, and $c \in \RR$,
    $$A + B = \bmat{a_1 + b_1 & a_2 + b_2 \\ 0 & a_3 + b_3} \in W$$
    $$cA = \bmat{ca_1 & ca_2 \\ 0 & ca_3} \in W$$
\end{itemize}
Therefore, $W \leq M_2(\RR)$ by the Subspace Test.
\end{proof}

\begin{example}
Let $V = M_2(\RR)$. Show that 
$$W = \{A \in M_2(\RR) \mid A \text{ is singular}\}$$
is \textbf{NOT} a subspace of $V$.
\end{example}

\begin{proof}
By counter example, let $A = \bmat{1 & 0 \\ 0 & 0}, B = \bmat{0 & 0 \\ 0 & 1}$,
$$A + B = \bmat{1 & 0 \\ 0 & 1} = I_2 \notin W$$
Therefore, $W \nleq V$ by the Subspace Test.
\end{proof}

\begin{example}
Let $V = P_1 = \{ax + b \mid a, b \in \RR\}$. Prove whether the following are subspaces of $P_1$ or not:
\begin{enumerate}[label=\alph*.]
    \item $W = \{ax - 2a \mid a \in \RR\}$
    \begin{proof}
        \leavevmode
        \begin{itemize}
            \item $0 = 0x = 0 \in W \rightarrow W \neq \emptyset$
            \item Let $f = ax - 2a, g = bx - 2b$, and $c \in \RR$. Then
            $$f + g = (a + b)x - 2(a + b) \in W$$
            $$cf = c(ax - 2a) = (ac)x - 2(ac) \in W$$
        \end{itemize}
    Therefore, $W \leq P_1$, by the Subspace Test.
    \end{proof}
    \item $W = \{ax + b \mid a + b = 1\}$
    \begin{proof}
    \leavevmode
    Take $f = x$ and $c = \dfrac{1}{2}$,
    $$cf = \frac{1}{2}x \notin W$$
    Therefore, $W \nleq P_1$, by the Subspace Test.
    \end{proof}
\end{enumerate}
\end{example}
\newpage
\section{Spanning Sets and Linear Independence}
Consider a vector space $V = \mathbb{R}^3$ and three vectors $\hat{\imath}, \hat{\jmath}$, and $\hat{k}$ where $\hat{\imath} = (1, 0, 0), \hat{\jmath} = (0, 1, 0), \hat{k} = (0, 0, 1) \in \mathbb{R}^3$. Then every vector $(a, b, c) \in \mathbb{R}^3$ can be expressed as 

$$(a, b, c) = (a, 0, 0) + (0, b, 0) + (0, 0, c)$$
$$=a\hat{\imath} + b\hat{\jmath} + c\hat{k}$$

Every vector in $\mathbb{R}^3$ can be expressed in terms of $\hat{\imath}, \hat{\jmath}$, and $\hat{k}$.

\begin{definition}
Let $V$ be a vector space and $v_1, v_2, \dots, v_k \in V$. A \textbf{linear combination} of $v_1, v_2, \dots, v_k$ is any element of the form
$$a_1v_1 + a_2v_2 + \cdots + a_kv_k = \sum_{i = 1}^k a_iv_i$$
where $a_i \in \mathbb{R}$ for all $i$.
\end{definition}

\begin{example} The following are examples of linear combinations:

\begin{enumerate}
\item $0v_1 + 0v_2 + \cdots + 0v_k = \vec{0}$ is a linear combination of $v_1, v_2, \dots, v_k$.

\item Every vector in $\mathbb{R}^3$ is a linear combination of $\hat{\imath}, \hat{\jmath}, \hat{k}$.

\item In the vector space $P_4$, $4x^4 - 2x^3 + 6x - 9 = 4x^4 + (-2)x^3 + 0x^2 + 6x + (-9)(1)$ is a linear combination of $1, x, x^2, x^3, x^4$.

\item Let $A = \bmat{1 & -3 & 0 \\ 0 & 2 & 2}$ and $x = \bmat{a \\ b \\ c}$. Then
\begin{align*}
Ax = \bmat{1 & -3 & 0 \\ 0 & 2 & 2}\bmat{a \\ b \\ c} = \bmat{a - 3b \\ 2b + 2c} \\
=\bmat{a \\ 0} + \bmat{-3b \\ 2b} + \bmat{0 \\ c} \\
= a\bmat{1 \\ 0} + b\bmat{-3 \\ 2} + c\bmat{0 \\ 2}
\end{align*}

$Ax$ is a linear combination of the columns of $A$. In general if $A \in M_{m \times n}(\mathbb{R})$ and $x \in M_{n \times 1}(\mathbb{R})$, then $Ax$ is a linear combination of the columns of $A$.

\item Let $v = (-3, -\frac{1}{2}, \frac{5}{2})$. Is $v$ a linear combination of $(1, 0, 0), (2, 1, 0)$, and $(0, 1, 1)$?

\textit{Solution.} We want to know if there exists scalars $a, b, c$ such that
$$\left(-3, -\frac{1}{2}, \frac{5}{2}\right) = a\left(1, 0, 0\right) + b(2, 1, 0) + c(0, 1, 1)$$
$$=(a + 2b, b + c, c)$$

Comparing coordinates,

\item In $V = \mathscr{F}(\mathbb{R})$, $\cosh \in V$ is a linear combination of $e^x$ and $e^{-x}$
$$\cosh x = \frac{e^x + e^{-x}}{2} = \frac{1}{2}e^x + \frac{1}{2}e^{-x}$$
\end{enumerate}


\end{example}

\begin{definition}
Let $S = \{v_1, v_2, \dots, v_k\}$ be a nonempty subset of a vector space $V$
\begin{itemize}
    \item The span of $S$, denoted by $\operatorname{span} S$, is the set of all linear combinations of elements of $S$, that is,
    $$\operatorname{span} S = \{a_1v_1 + a_2v_2 + \cdots + a_kv_k \mid a_i \in \mathbb{R}\,\, \forall i\}$$
    \item If $\operatorname{span} S = V$, we say that $S$ is a \underline{spanning set} of $V$ or $V$ is spanned by $S$.\footnote{Every vector in $V$ is a linear combination of elements in $S$.}
\end{itemize}
\end{definition}

\begin{example}
Let $V = \mathbb{R}^2$ and $S = \{(1, 0, 0), (0, 1, 0), (0, 0, 1)\}$, then,
$$\operatorname{span}S = \{a(1,0, 0) + b(0, 1, 0) + c(0, 0, 1) \mid a, b, c \in \mathbb{R}\}$$
$$=\{(a, b, c) \mid a, b, c \in \mathbb{R}\}$$
Therefore, $S$ spans $\mathbb{R}^3$.
\end{example}

\begin{example}
Let $V = P_n$ and $S = \{1, x, x^2, \dots, x^n\}$, then,
$$\operatorname{span} S = \{a_0 + a_1x + a_2x^2 + \cdots + a_nx^n \mid a_i \in \RR\,\, \forall i\} = P_n$$
Therefore, $S$ spans $P_n$.
\end{example}

\begin{example}
Let $V = \CC = \{a + bi \mid a, b \in \RR\,\,\text{and}\,\,i^2 = -1\}$ and $S = \{1, i\}$, then,
$$\operatorname{span}S = \{a + bi \mid a, b \in \RR\,\,\text{and}\,\,i^2 = -1\} = \CC$$
Therefore, $S$ spans $\CC$.
\end{example}

\begin{example}
Let $V = M_2(\RR)$ and $S = \left\{\bmat{1 & 0 \\ 0 & 0}, \bmat{ 0 & 1 \\ 0 & 0}, \bmat{0 & 0 \\ 1 & 0}, \bmat{0 & 0 \\ 0 & 1}\right\}$, then,
$$\operatorname{span}S = \left\{\bmat{a & b \\ c & d} \suchthat a,b, c, d \in \RR \right\} = M_2(\RR)$$
Therefore $S$ spans $M_2(\RR)$.
\end{example}

\begin{example}
Let $S = \{(2, 3), (1, 2)\}$. Does this span $\RR^2$?
\end{example}

\textit{Solution.} Let $(a, b) \in \RR^2$. We want to know if there exists $x, y \in \RR$ such that,
$$(a, b) = x(2, 3) + y(1, 2)$$
$$=(2x, 3x) + (y, 2y)$$
$$=(2x + y, 3x + 2y)$$
\begin{align*}
\begin{cases}
    2x + y = a \\
    3x + 2y = b 
\end{cases}
\quad
\longleftrightarrow
\quad
\bmat{2 & 1 \\ 3 & 2}\bmat{x \\ y} = \bmat{a \\ b}
\end{align*}
The coefficient matrix $A$ of the system is nonsingular, since $\det A \neq 0$, therefore the system has a solution.
$$\bmat{x \\ y} = \bmat{2 & 1 \\ 3 & 2}\inv \bmat{a \\ b}$$
$$= \bmat{2 & -1 \\ 3 & -2}\bmat{a \\ b} = \bmat{2a - b \\ 3a - 2b}$$
Therefore, $S$ spans $\RR^2$.

\newpage

\begin{example}
Let $S = \{x^2 + 2x + 1, x^2 + 2\} \subseteq P_2$. Does $S$ span $P_2$?
\end{example}

\textit{Solution.} Let $ax^2 + bx + c \in P_2$. We want to know if there exists $D, E \in \RR$ such that
$$ax^2 + bx + c = D(x^2 + 2x + 1) + E(x^2 + 2)$$
$$=x^2(D + E) + x(2D) + (D + 2E)$$
\begin{align*}
\begin{cases}
    D + E = a \\ 
    2D = b \\
    D + 2E = c
\end{cases}
\quad
\longleftrightarrow
\quad 
\begin{amatrix}{2}
    1 & 0 & a \\
    2 & 0 & b \\
    1 & 2 & c
\end{amatrix}
\xrightarrow{EROs}
\begin{amatrix}{2}
    1 & 1 & a \\
    0 & 1 & -a+ c \\
    0 & 0 & -4a + b + 2c \\
\end{amatrix}
\end{align*}
Since there is a restriction placed on the last row, such that $-4a + b + 2c$ must be equal to 0, therefore $S$ does not span $P_2$.

\vspace{0.2cm}

\begin{theorem}
If $S = \{v_1, v_2, \dots, v_k\}$ is a nonempty subset of $V$, then $\operatorname{span}S \leq V$.
\end{theorem}

\begin{remark}Shown below are some remarks regarding the span of sets:
\begin{itemize}
    \item If $S \subseteq \operatorname{span} S$ and $\operatorname{span}S$ is the smallest subspace of $V$ containing $S$.
    \item $\operatorname{span}(\emptyset) = \{\vec{0}\}$
\end{itemize}
\end{remark}

\begin{definition}
Let $S = \{v_1, v_2, \dots, v_n\}$ be a nonempty subset of a vector space $V$. Then $S$ is said to be \textbf{linearly independent} if whenever
$$a_1v_1 + a_2v_2 + \cdots + a_nv_n = \vec{0}$$
then $a_1 = a_2 = \cdots = a_n = 0$. Otherwise, $S$ is said to be \textbf{linearly dependent}.
\end{definition}

\begin{remark}
The following are some remarks on linear dependence:
\begin{enumerate}
    \item $S$ is linearly dependent iff there exists $a_1, a_2, \dots, a_n \in \RR$ and not all zero, such that,
    $$\sum_{i=1}^na_iv_i = \vec{0}$$
    \item $\emptyset$ is taken to be linearly independent.
\end{enumerate}
\end{remark}

\begin{example}\label{example:6.5}
Determine if the set $S = \{(1, 0, 0), (0, 1, 0), (0, 0, 1)\} \subseteq \mathbb{R}^3$ is linearly independent.

\textit{Solution.} Form the equation,
$$a(1, 0, 0) + b(0, 1, 0) + c(0, 0, 1) = (0, 0, 0)$$
$$(a, b, c) = (0, 0, 0)$$
Because $a = b = c = 0$, therefore $S$ is linearly independent.

In general, $S = \{e_1, e_2, \dots, e_n\}$ is a linear independent subset of $\RR^n$.
\end{example}

\begin{example}\label{example:6.6}
Determine whether the set $S = \left\{\bmat{1 & 0 \\ 0 & 0}, \bmat{0 & 1 \\ 0 & 0}, \bmat{0 & 0 \\ 1 & 0}, \bmat{0 & 0 \\ 0 & 1}\right\} \subseteq M_2(\RR)$ is linearly independent.

\textit{Solution.} Form the equation,
$$a\bmat{1 & 0 \\ 0 & 0} + b\bmat{0 & 1 \\ 0 & 0} + c\bmat{0 & 0 \\ 1 & 0} + d\bmat{0 & 0 \\ 0 & 1} = \bmat{0 & 0 \\ 0 & 0}$$
$$=\bmat{a & 0 \\ 0 & 0} + \bmat{0 & b \\ 0 & 0} + \bmat{0 & 0 \\ c & 0} + \bmat{0 & 0 \\ 0 & d} = \bmat{0 & 0 \\ 0 & 0}$$
$$=\bmat{a & b \\ c & d} = \bmat{0 & 0 \\ 0 & 0}$$
where $a, b, c, d \in \RR$. Since $a = b = c = d = 0$, therefore, $S$ is linearly independent.
\end{example}

\lline

In general, $S = \{E_{k\ell} \mid k = 1, 2, \dots, m, \quad \ell = 1, 2, \dots, n\}$ is a linear independent subset of $M_{m \times n}(\RR)$. As an exercise, determine if the set $S$ is a spanning set in \hyperref[example:6.5]{Example 6.5} and \hyperref[example:6.6]{Example 6.6}.

\begin{example}
Determine if the set $S = \{e^x, e^{-x}\} \subseteq \mathscr{F}(\RR)$.

\textit{Solution.} Form the equation,
$$ae^x + be^{-x} = 0$$
where $a, b \in \RR$. By differentiating both sides we obtain,
$$ae^x -be^{-x} = 0$$
Adding the first and second equation together,
$$ae^x + be^{-x} + ae^x - be^{-x} = 0$$
$$2ae^x = 0 \Longrightarrow a = 0$$
Since $a = b = 0$, therefore, $S$ is linearly independent.
\end{example}

\begin{remark}
Let $V$ be a vector space and $S, T \leq V$ be finite:
\begin{enumerate}
    \item If $\vec{0} \in S$, then $S$ is linearly dependent.
    \item If $\vec{0} + v \in V$, then $\{v\}$ is linearly independent.
    \item Suppose $S \subseteq T$
    \begin{enumerate}
        \item $\operatorname{span}S \subseteq \operatorname{span}T$
        \item If $T$ is linearly independent then $S$ is linearly independent
        \item If $S$ is linearly dependent then $T$ is linearly dependent
    \end{enumerate}
    \item $S$ is linearly dependent iff some vector in $S$ is a linear combination of the other vectors in $S$. In particular, $S = \{u, v\}\quad(u, v \ne 0)$ is linearly dependent iff $v = cu$ for some $0 \ne c \in \RR$. In this case, $u$ and $v$ are said to be parallel.
\end{enumerate}
\end{remark}
% Might need to add more
\newpage
\section{Basis and Dimension of a Vector Space}
\begin{definition}
A subset $B$ of a vector space $V$ is called a basis for $V$ if $\operatorname{span}B = V$ and $B$ is linearly independent.
\end{definition}


\begin{example}
\leavevmode
\begin{enumerate}
    \item $\emptyset$ is a basis for $\{\vec{0}\}$
    \item $B = \{e_1, e_2, \dots, e_n\} \subseteq \RR^n$ is a basis for $\RR^n$
    \item $B = \{E_{k\ell} \mid k = 1, 2, \dots, m \quad l = 1, 2, \dots, n\}$ is a basis for $\matmn$ 
    \item $B = \{1, x, x^2, \dots, x^n\}$ is a basis for $P_n$
    \item $B = \{1, i\}$ is a basis for $\CC$.
\end{enumerate}
\end{example}

\begin{remark}
The bases in Example 2 - 5 as shown above are known as the \textbf{standard bases} of their respective spaces.
\end{remark}

\begin{example}
Is $B = \{(2, 3), (1, 2)\}$ a basis for $\RR^2$?
\end{example}

\textit{Solution.} Recall that a set $B$ is a basis if $\operatorname{span} B = V$ and $B$ is linearly independent. Let us verify whether these conditions hold true:
$$x(2, 3) + y(1, 2) = (a, b)$$
$$(2x + y, 3x + 2y) = (a, b)$$

To verify, whether $B$ is linearly independent,
$$(2x + y, 3x + 2y) = (0, 0)$$
$$\begin{amatrix}{2}
    2 & 1 & 0 \\
    3 & 2 & 0
\end{amatrix} \xrightarrow{EROs}
\begin{amatrix}{2}
    2 & 1 & 0 \\ 
    1 & 0 & 0 \\
\end{amatrix}$$
Because $\det B \neq 0$ it is nonsingular and its solution is the trivial solution. Therefore, $B$ is linearly independent and is a basis for $\RR^2$.

\begin{example}
Is $B = \{x^2 + x, x^2 + 1, 1\}$ a basis for $P_2$?
\end{example}

\textit{Solution.} Let us first verify if $B$ is a spanning set. Let $ax^2 + bx + c \in P_2$, and $D, E, F \in \RR$,
$$ax^2 + bx + c = D(x^2 + x) + E(x^2 + 1) + F(1)$$
$$=x^2(D + E) + x(D) + (E + F)$$
$$\begin{cases}
    D + E = a \\
    D = b \\
    E + F = c
\end{cases}
\quad \longleftrightarrow \quad
\begin{amatrix}{3}
    1 & 1 & 0 & a \\
    1 & 0 & 0 & b \\
    0 & 1 & 1 & c \\
\end{amatrix}
$$

Since the system of equations has a solution with the form $(a, a - b, c - a + b)$. Therefore, it spans $P_2$. To verify that it is linearly independent,
$$0 = x^2(D + E) +x(D) + (E + F)$$
Since the system of equations forces the other variables to be 0, thus, $B$ is linearly independent. Therefore, $B$ is a basis for $P_2$.\footnote{An alternative solution would be to use determinants.}


\newpage
\begin{theorem}
If $B = \{v_1, v_2, \dots, v_k\}$ is a basis for a vector space $V$, then every vector in $V$ can be expressed in exactly one way as a linear combination of vectors in $B$.
$$v = a_1v_1 + a_2v_2 + \cdots + a_kv_k$$
\end{theorem}

\begin{theorem}
Let $V$ be a vector space having a finite spanning set $S$. Then $S$ has a subset $B$ that is a basis for $V$. To find a basis for $V$ with finite spanning set $S = \{v_1, v_2, \dots, v_k\}$
\begin{enumerate}
    \item Form $\vec{0} = x_1v_1 + x_2v_2 + \cdots + x_kv_k$ where each $x_i \in \RR$ is unknown.
    \item Use EROs to bring the coefficient matrix to REF or RREF.
    \item The vectors corresponding to columns containing the leading entries forms a basis for $V$.
\end{enumerate}
\end{theorem}

\lline

\begin{example}
Find a basis for the given subspace
$$W:=\spn\{x^2, x + 4, 2x + 7, 1\} \leq P_2$$
\textit{Solution.}

A spanning set for $W$ is $S = \{x^2, x + 4, 2x + 7, 1\}$
$$\vec{0} = Ax^2 + B(x + 4) + C(2x + 7) + D$$
$$\vec{0} = x^2(A) + x(B + 2C) + (4B + 7C + D)$$
Forming the augmented matrix,
$$\begin{amatrix}{4}
    1 & 0 & 0 & 0 & 0 \\
    0 & 1 & 2 & 0 & 0 \\
    0 & 4 & 7 & 1 & 0
\end{amatrix} \xrightarrow{\text{RREF}}
\begin{amatrix}{4}
    1 & 0 & 0 &0 & 0 \\
    0 & 1 & 0 & 2 & 0 \\ 
    0 & 0 & 1 & -1 & 0 \\
\end{amatrix}
$$
And since the fourth column has no leading entries, we can remove it. Thus, $B = \{x^2, x + 4, 2x + 7\}$ is a basis for $W$.
\end{example}

\begin{example}
Find a basis for the given subspace $W = \spn\{(1, -1, 3, 2), (0, -1, 2, 1), (2, 1, 0, 1)\} \subseteq \RR^2$.

\textit{Solution.}
$$a(1, -1, 3, 2) + b(0, -1, 2, 1) + c(2, 1, 0, 1) = (0, 0, 0, 0)$$
$$(a + 2c, -a - b + c, 3a + 2b, 2a + b + c) = (0, 0, 0,0)$$
$$\begin{amatrix}{3}
1 & 0 & 2 & 0 \\
-1 & -1 & 1 & 0 \\
3 & 2 & 0 & 0 \\
2 & 1 & 1 & 0
\end{amatrix} \xrightarrow{\text{RREF}}
\begin{amatrix}{3}
    1 & 0 & 2 & 0 \\
    0 & 1 & -3 & 0 \\
    0 & 0 & 0 & 0 \\
    0 & 0 & 0 & 0 \\
\end{amatrix}
$$
Therefore, $\{(1, -1, 3, 2), (0, -1, 2, 1)\}$ is a basis for $W$.
\end{example}

\newpage
\begin{example}
Find a basis for the given subspace $W = \left\{\bmat{a & b \\ c & d} \in M_2(\RR) \suchthat b + 2c - d = 0\right\}$.

\textit{Solution.}

Let us first find a spanning set,
\begin{align*}
    A = \bmat{a & b \\ c & d} \in W &\longleftrightarrow  b+ 2c - d = 0 \\
    &\longleftrightarrow \bmat{a & b \\ c & b + 2c} \\
    &\longleftrightarrow \bmat{a & 0 \\ 0 & 0 } + \bmat{0 & b \\ 0 & b} + \bmat{0 & 0 \\ c & 2c} \\
    &\longleftrightarrow a\bmat{1 & 0 \\ 0 & 0} + b\bmat{0 & 1 \\ 0 & 1} + c\bmat{0 & 0 \\ 1 & 2}
\end{align*}
$$W = \spn\underbrace{\left\{\bmat{1 & 0 \\ 0 & 0}, \bmat{0 & 1 \\ 0 & 1}, \bmat{0 & 0 \\ 1 & 2}\right\}}_B$$
Thus, $B$ spans $W$. Form,
$$a\bmat{1 & 0 \\ 0 & 0} + b\bmat{0 & 1\\ 0 & 1} + c\bmat{0 & 0 \\ 1 & 2} = \bmat{0 & 0 \\ 0 & 0}, \quad a, b, c \in \RR$$
$$\bmat{a & b \\ c & b + 2c} = \bmat{0 & 0 \\ 0 & 0}$$
$$\Longrightarrow \begin{cases}
    a = 0 \\ 
    b = 0 \\
    c = 0 \\
    b + 2c = 0
\end{cases}$$
And since $a = b = c = 0$, $B$ is linearly independent. Therefore, $B$ is a basis for $W$.
\end{example}

\begin{example}
Find a basis for $W = \{ax^3 + bx^2 + cx + d \in P_3 \mid a + c = 0 \text{ and } b + d = 0\}$.

\textit{Solution.}

Let $f(x) = ax^3 + bx^2 + cx + d \in W$ and since there are restrictions, it must follow that $c = -a$ and $d = -b$,
$$ax^3 + bx^2 -ax -b = a(x^3 - x) + b(x^2 - 1)$$
$$f(x) \in \spn\{x^3 - x, x^2 - 1\}$$
$$\therefore W = \spn\underbrace{\{x^3 - x, x^2 - 1\}}_B$$
Thus, $B$ spans $W$. And since $x^2 - 1$ is not a scalar multiple of $x^3 - x$, $B$ is linearly independent. Therefore, $B$ is a basis for $W$.
\end{example}

\begin{definition}
A vector space $V$ with a finite spanning set is said to be \textbf{finite dimensional}. Otherwise, $V$ is said to be \textbf{infinite dimensional}.
\end{definition}

\newpage
\begin{remark}The following are some remarks on basis:
\begin{enumerate}
    \item Every finite-dimensional vector space has a finite basis.
    \item Every infinite-dimensional vector space has an infinite basis.\\
    \underline{\textbf{Example:}}
    \leavevmode
    \begin{itemize}
        \item The vector space $\RR^n, M_{m \times n}(\RR), P_n, \CC$ are all finite-dimensional.
        \item $\mathscr{F}(\RR)$ is infinite-dimensional.
        \item $\RR[x]$ the set of all polynomials in the variable $x$ having real coefficients is an infinite-dimensional vector space under addition and scalar multiplication of polynomials. ($B = \{x^n \mid 0 \leq n \in \ZZ\}$ is an infinite basis for $\RR[x]$).
    \end{itemize}
\end{enumerate}
\end{remark}

\begin{theorem}[Basis Completion Theorem]
Let $V$ be a finite-dimensional vector space and $S \subseteq V$ be a linearly independent subset of $V$. Then $V$ has a basis $B$ containing $S$.

To extend a linearly independent set $S$ to a basis for $V$
\begin{enumerate}
    \item Let $C$ be a known basis (standard basis) for $V$. Then $T := S \cup C$ is a spanning set for $V$.
    \item Use the preceding algorithm to get a subset $B$ of $T$ that is a basis for $V$, Then $B$ contains $S$.
\end{enumerate}
\end{theorem}

\begin{example}
Find a basis for $M_{3 \times 1}(\RR)$ containing $v_1 = \bmat{2 \\ 1 \\ 0}$ and $v_2 = \bmat{1 \\ -1 \\ 1}$.
\end{example}

\textit{Solution.} Verify that $S = \{v_1, v_2\}$ is linearly independent.

Let $C = \left\{\bmat{1 \\ 0 \\ 0}, \bmat{0 \\ 1 \\ 0}, \bmat{0 \\ 0 \\ 1}\right\}$ be the standard basis for $M_{3 \times 1}(\RR)$.

Form
$$T = \left\{\bmat{2 \\ 1 \\ 0}, \bmat{1 \\ -1 \\ 1}, \bmat{1 \\ 0 \\ 0}, \bmat{0 \\ 1 \\ 0}, \bmat{0 \\ 0 \\ 1}\right\}$$

Using the previous algorithm, we obtain the following coefficient matrix,
$$\bmat{2 & 1 & 1 & 0 & 0 \\ 1 & -1 & 0 & 1 & 0 \\ 0 & 1 & 0 & 0 & 1} \xrightarrow{\text{RREF}} \bmat{1 & 0 & 1/3 & 1/3 & 0 \\ 0 & 1 & 1/3 & -2/3 & 0 \\ 0 & 0 & 0 & 0 & 1}$$

Since there are no pivots for columns 3 and 4. We can delete them. Therefore, $B = \left\{\bmat{2 \\ 1 \\ 0}, \bmat{0 \\ 1 \\ 0}, \bmat{0 \\ 0 \\ 1}\right\}$ is a basis for $M_{3 \times 1}(\RR)$ containing $v_1$ and $v_2$.

\begin{theorem}[Steinite Replacement Theorem]
Let $V$ be a finite-dimensional vector space with finite spanning set $S$. If $T$ is a linearly independent subset of $V$, then $|T| \leq |S|$.
\end{theorem}

\begin{corollary}
Any two bases of a finite-dimensional vector space $V$ has the same number of elements.
\end{corollary}

\begin{definition}
The number of elements of any basis of a finite-dimensional vector space $V$ is called the dimension of $V$, denoted by $\dim V$. If $n = \dim V$, then $V$ is said to be $n$-dimensional.
\end{definition}

\begin{remark}
If $V$ is an infinite-dimensional vector space, then $\dim V = +\infty$.
\end{remark}
\newpage
\section{The Fundamental Subspaces and Rank of a Matrix}

\begin{definition}
Let $A \in M_{m \times n}(\RR)$ with rows $R_1, R_2, \dots, R_m$ and columns $C_1, C_2, \dots, C_n$.
\begin{enumerate}
    \item The \textbf{row space} of $A$ is $R(A) := \operatorname{span}\{R_1, R_2, \dots, R_m\} \leq M_{1 \times n}(\RR)$
    \item The \textbf{column space} of $A$ is $C(A) := \operatorname{span}\{C_1, C_2, \dots, C_n\} \leq M_{m \times 1}(\RR)$
    \item The \textbf{row rank} (respectively \textbf{column rank}) of $A$ is $\dim R(A)$ (respectively $\dim C(A)$)
\end{enumerate}
\end{definition}
\leavevmode \\
To find bases for $R(A)$ and $C(A)$:
\begin{enumerate}
    \item Use elementary row operations to bring $A$ to a matrix $U$ in REF.
    \item Basis for $C(A) = $ columns of $A$ corresponding to columns of $U$ containing leading entries.
    \item Basis for $R(A) = $ nonzero rows of the REF matrix $U$.
\end{enumerate}

\begin{theorem}
The row rank and column rank of any $A \in M_{m \times n}(\RR)$ are equal.
\end{theorem}

\begin{definition}
The rank of $A \in M_{m \times n}(\RR)$, denoted by $\rank A$, is the common value of the row rank and column ranks of $A$.
\end{definition}

\begin{definition}
Let $A \in M_{m \times n}(\RR)$. The null space of $A$, denoted by $N(A)$, is the solution set of the system $Ax = 0_{m \times 1}$. That is,
$$N(A) = \{x \in M_{n \times 1}(\RR) \mid Ax = 0_{m \times 1}\}$$
\end{definition}

\begin{remark}
If $A \in M_{m \times n}(\RR)$, then $N(A) \leq M_{n \times 1}(\RR)$. The \textbf{nullity} of $A$, denoted by $\upsilon(A)$, is $\dim N(A)$.
\end{remark}

To find $N(A)$, solve for $x$ in the system $Ax = 0_{m \times 1}$.

\begin{example}
For each matrix $A$, find a basis for $N(A)$. Identify the nullity of $A$.
$$A = \bmat{1 & 0 \\ 4 & -2 \\ 2 & 3}$$
\end{example}

\textit{Solution.} 

By reducing matrix $A$ to RREF, we get
$$\bmat{1 & 0 \\ 0 & 1 \\ 0 & 0} \longrightarrow 
\begin{cases}
x_1 = 0 \\
x_2 = 0
\end{cases}
$$
Therefore, $N(A) = \left\{\bmat{0 \\ 0}\right\}$ and $\upsilon(A) = 0$.

\begin{example}
Find a basis for the given system of equations
$$
\begin{cases}
x_1 + 2x_3 + x_5 = 0 \\
x_2 + x_3 + x_5 = 0 \\
x_4 - x_5 = 0
\end{cases}
$$
\end{example}

\begin{theorem}[Rank-Nullity Theorem]
If $A \in M_{m \times n}(\RR)$, then $n = \rank A + \upsilon(A)$,
\end{theorem}
\newpage
\section{Inner Product Spaces}
\begin{definition}
An inner product on a vector space $V$ is a function $\langle \cdot, \cdot \rangle$ that assigns two vectors $u, v \in V$ to a scalar $\langle u, v \rangle \in \RR$ satisfying the following for all $u, v, w \in V$ and for all $c \in \RR$

IP1. $\langle u, v \rangle = \langle v, u \rangle$ (symmetric)

IP2. $\langle u, u \rangle \geq 0$ (nonnegativity)

IP3. $\langle u, u \rangle = 0$ iff $u = \vec{0}$ (definite)

IP4. $\langle u + v, w \rangle = \langle u, w \rangle + \langle v, w \rangle$  (additive on the first component)

IP5. $\langle cu, w \rangle = c\langle u, w \rangle$ (homogeneous on the first component)
\end{definition}

\begin{remark}
A vector space $V$ with an inner product is called an inner product space.\footnote{Also known as a Euclidean space.}
\end{remark}

\begin{definition}
Let $V$ be an inner product space and $u, v \in V$
\begin{enumerate}
    \item The \textbf{norm/magnitude} of $u$ is $\lVert u \rVert = \sqrt{\langle u, u \rangle}$
    \item $u$ is called the \textbf{unit vector} if $\lVert u \rVert = 1$
    \item The \textbf{distance} between $u$ and $v$ is $d(u, v) = \lVert u - v \rVert$
\end{enumerate}
\end{definition}

\begin{remark}
Let $V$ be an inner product space, $u \in V$, $c \in \RR$
\begin{enumerate}
    \item $\lVert u \rVert^2 = \langle u, u \rangle$
    \item $\lVert u \rVert \geq 0$
    \item $\lVert \vec{0} \rVert = \sqrt{\langle \vec{0}, \vec{0} \rangle} = 0$
    \item $\lVert u \rVert = 0$ iff $u = \vec{0}$
    \item $\lVert cu \rVert = \sqrt{\langle cu, cu \rangle} = \lvert c \rvert \lVert u \rVert$
    \item If $u + \vec{0}$, then $\dfrac{u}{\lVert u \rVert}$ is a unit vector, called the normalization of $u$
    \item $\lVert u + v \rVert^2 = \langle u + v, u + v \rangle = \langle u, u \rangle + \langle u, v \rangle + \langle v, u \rangle + \langle v, v \rangle =\lVert u \rVert^2 + 2\langle u, v \rangle + \lVert v \rVert^2$ \\
    $\lVert u - v \rVert^2 = \lVert u \rVert^2 - 2 \langle u, v \rangle + \lVert v \rVert^2$
\end{enumerate}
\end{remark}

\begin{theorem}[Cauchy-Schwarz Inequality]
Let $V$ be an inner product space and $u, v \in V$. Then $\lvert \langle u, v \rangle \rvert \leq \lVert u \rVert \lVert v \rVert$.
\end{theorem}

\begin{remark}
Let $u$ and $v$ be nonzero vectors in an inner product space $V$
\begin{enumerate}
    \item $$-\lVert u \rVert \lVert v \rVert \leq \langle u, v \rangle \leq \lVert u \rVert \lVert v \rVert$$
    $$-1 \leq \langle u, v \rangle \leq 1$$
    \item The angle $\theta$ between $u$ and $v$ is defined as 
    $$\theta = \cos^{\inv}\left(\frac{\langle u, v \rangle}{\lVert u \rVert \lVert v \rVert}\right)$$
\end{enumerate}
\end{remark}

\newpage
\begin{theorem}[Properties of the Norm]
Let $V$ be an inner product space, $u, v \in V$, and $c \in \RR$.
\begin{enumerate}
    \item $\lVert u \rVert \geq 0$
    \item $\lVert u \rVert = 0$ iff $u = \vec{0}$
    \item $\lVert cu \rVert = \lvert c \rvert \lVert u \rVert$
    \item $\lVert u + v \rVert \leq \lVert u \rVert + \lVert v \rVert$ (Triangle inequality)
    \begin{proof}
    Recall that, $\lVert u + v \rVert^2$ can be expressed as $\langle u + v, u + v \rangle$,
    $$\langle u + v, u + v \rangle = \lVert u \rVert^2 + 2 \lVert u \rVert \lVert v \rVert + \lVert v \rVert^2 \leq (\lVert u \rVert + \lVert v \rVert)^2$$
    \end{proof} 
\end{enumerate}
\end{theorem}

\begin{remark}[Properties of Distance]
Let $V$ be an inner product space and $u, v, w \in V$. Then
\begin{enumerate}
    \item $d(u, v) \geq 0$
    \item $d(u, v) = 0$ iff $u = v$
    \item $d(u, v) = d(v, u)$
    \item $d(u, w) \leq d(u, v) + d(v, w)$
\end{enumerate}
\end{remark}

\begin{definition}
Let $V$ be an inner product space, $u, v \in V$ and $S \subseteq V$ be nonempty. 
\begin{enumerate}
    \item $u$ is \textbf{orthogonal} to $v$, denoted by $u \perp v$, iff $\langle u, v \rangle = 0$
    \item $S$ is an \textbf{orthogonal set} if for all $u, v \in V$, $u \neq v$, and $u \perp v$
    \item $S$ is an \textbf{orthonormal set} if $S$ is orthogonal and every vector in $S$ is a unit vector.
\end{enumerate}
\end{definition}

\begin{example}
Let $V = \RR^3$ under the standard inner product and 
$$S = \{(0, 1, 0), (1, 0, 1), (1, 0, -1)\}, \quad T = \left\{(0, 1, 0), \left(-\frac{4}{5}, 0, \frac{3}{5}\right), \left(\frac{3}{5}, 0, \frac{4}{5}\right)\right\}$$

If we take all 2-combination of $S$, then we can show that every pair is orthogonal to each other, but not orthonormal. However, in $T$, we can show that every pair is orthogonal to each other and each vector is a unit vector, thus, it is orthonormal.
\end{example}

\begin{example}
Let $V = C[-\pi, \pi]$, then $S = \{\sin x, \cos x\}$ is orthogonal, but not orthonormal.\footnote{Prove this as an exercise}
\end{example}

\begin{remark}
Let $V$ be an inner product space:
\begin{enumerate}
    \item For all $u \in V$, $\vec{0} \perp u$
    \item If $u \in V$ such that $u \perp s$, for all $s \in V$, then $u = \vec{0}$
    \item If $u \perp v$ and $w \perp v$, then $au + bw \perp v$, for all $a, b \in \RR$.
    \item $S = \{v_1, v_2, \dots, v_n\}$ is an orthogonal set iff
    $$\langle v_i, v_j \rangle = \begin{cases}
        0 & \text{if $i = j$} \\
        1 & \text{if $i \neq j$}
    \end{cases} = \delta_{ij}$$
\end{enumerate}
\end{remark}

\begin{theorem}[Pythagorean Theorem]
Let $V$ be an inner product space and $u, v \in V$. Then $u \perp v$ iff $\lVert u + v \rVert^2 = \lVert u \rVert ^2 + \lVert v \rVert^2$.
\end{theorem}

\newpage
\begin{definition}
Let $V$ be an inner product space, $u \in V$, and $S \subseteq V$
\begin{enumerate}
    \item $u$ is orthogonal to $S$, denoted by $u \perp S$, if $u \perp s$ for all $s \in S$
    \item The orthogonal complement of $S$, denoted by $S^\perp$ (S-perp) is
    $$S^\perp = \{u \in V \mid u \perp s, \hspace{0.2cm} \forall s \in S\}$$
\end{enumerate}
\end{definition}

\begin{remark}
Let $V$ be an inner product space
\begin{enumerate}
    \item 
    \begin{enumerate}
        \item $\{0\}^\perp = \{u \in V \mid \vec{0} \perp u\} = V$
        \item $V^\perp = \{u \in V \mid u \perp s, \hspace{0.2cm} \forall s \in S\} = \{\vec{0}\}$
    \end{enumerate}
    \item If $S \subseteq V$, then $S^\perp = (\operatorname{span} S)^\perp$. In particular if $W = \operatorname{span} S$, then $W^\perp = S^\perp$
\end{enumerate}
\end{remark}

\begin{theorem}
If $S$ is a subset of an inner product space, then $S^\perp \leq V$.
\end{theorem}

\begin{example}
Find the orthogonal complement of $W_1 = \{(c - 6d, d, c) \mid c, d \in \RR\} \subseteq \RR^3$ under the standard inner product.
\end{example}

\textit{Solution.} 

$$W_1 = \{c(1, 0, 1) + d(-6, 1, 0) \mid c, d \in \RR\} = \operatorname{span}\{(1, 0, 1), (-6, 1, 0)\}$$

Thus, $W_1^\perp = S^\perp$ where $S = \{(1, 0, 1), (-6, 1, 0)\}$. Now recall that for a vector to be in the orthogonal complement,
$$(x, y, z) \in S^\perp \longleftrightarrow \begin{cases}
    \langle (x, y, z), (1, 0, 1)\rangle = 0 \\
    \langle (x, y, z), (-6, 1, 0)\rangle = 0
\end{cases}$$

Forming the system of equations,
$$\begin{cases}
    x + z = 0 \\
    -6x + y = 0
\end{cases} \quad \longleftrightarrow \quad \bmat{1 & 0 & 0 \\ -6 & 1 & 0} \xrightarrow{EROs} \bmat{1 & 0 & 1 \\ 0 & 1 & 6}$$

Let $z = t$, then $x = -t$ and $y = -6t$,
$$W_1^\perp = S^\perp = \{(-t, -6t, t) \mid t \in \RR\} = \{t(-1, -6, 1) \mid t \in \RR\}$$
$$=\boxed{\operatorname{span}\{(-1, 6, 1)\}}$$

\begin{theorem}
A finite orthogonal set $S = \{v_1, v_2, \dots, v_n\}$ of nonzero vectors in an inner product space $V$ is linearly independent.
\end{theorem}

\begin{proof}
Suppose that,
$$c_1v_1 + c_2v_2 + \cdots + c_nv_n = \vec{0}$$
To demonstrate that $S = \{v_1, v_2, \dots, v_n\}$ is linearly independent, we must show that $c_1 = c_2 = \cdots = c_n = 0$. Recall that,
$$\langle\vec0, v_i\rangle = 0 \Longrightarrow \langle c_1v_1 + c_2v_2 + \cdots + c_nv_n, v_i \rangle$$
$$c_1\langle v_1, v_i \rangle +c_2\langle v_2, v_i \rangle + \cdots + c_n\langle v_n, v_i\rangle$$
From the orthogonality of $S$ it follows that $\langle v_i, v_j \rangle = 0$ if $i \neq j$, and it is nonzero if $i = j$. Thus, the equation reduces to
$$c_i\langle v_i, v_i\rangle = 0$$
Since $S$ does not contain nonzero vectors, then this implies that each $c_i$ in the equation is equal to zero.
\end{proof}


\begin{definition}
Let $V$ be a finite-dimensional inner product space. A basis $B$ for $V$ is called an \textbf{orthogonal basis} (resp. orthonormal basis) for $V$ if $B$ is an orthogonal (resp. orthonormal) set.
\end{definition}

\begin{theorem}
Let $V$ be an inner product space and $B = \{v_1, v_2, \dots, v_n\}$ be an orthogonal basis for $V$. Then every $v \in V$ can be written as a linear combination
$$v = c_1v_1 + c_2v_2 + \cdots + c_nv_n$$
where 
$$c_i = \frac{\langle v, v_i\rangle}{\langle v_i, v_i \rangle}$$
\end{theorem}

\begin{example}
Let $B = \{\bmat{-1 & 1 & 1 & 1}^{\mathsf{T}}, \bmat{1 & -1 & 1 & 1}^{\mathsf{T}}, \bmat{1 &1 & -1 & 1}^{\mathsf{T}}, \bmat{1 & 1 & 1 & -1}^{\mathsf{T}}\}$.

\begin{enumerate}
    \item Show that $B$ is an orthogonal basis of $M_{4 \times 1}(\RR)$ under the standard inner product.

    \textit{Solution.}
    \begin{align*}
    \langle \bmat{-1 & 1 & 1 & 1}^{\mathsf{T}}, \bmat{1 & -1 & 1 & 1}^{\mathsf{T}} \rangle = 0, \quad \quad \langle\bmat{-1 & 1 & 1 & 1}, \bmat{1 & 1 & -1 & 1}^{\mathsf{T}}\rangle = 0 \\
    \langle\bmat{-1 & 1 & 1 & 1}^{\mathsf{T}}, \bmat{1 & 1 & 1 & -1}^{\mathsf{T}}\rangle = 0, \quad \quad \langle\bmat{1 & -1 & 1 & 1}, \bmat{1 & 1 & -1 & 1}^{\mathsf{T}}\rangle = 0 \\
    \langle\bmat{1 & -1 & 1 & 1}^{\mathsf{T}}, \bmat{1 & 1 & 1 & -1}^{\mathsf{T}}\rangle = 0, \quad \quad \langle\bmat{1 & 1 & -1 & 1}^{\mathsf{T}}, \bmat{1 & 1 & 1 & -1}^{\mathsf{T}}\rangle = 0
    \end{align*}
    
    \item Use (1) to solve for $c$ in
    $$\begin{cases}
        -a + b + c + d = 1\\
        a - b + c + d = 8 \\
        a + b - c + d = -7 \\
        a + b + c - d = 2
    \end{cases}$$

    \textit{Solution.} The given system is equivalent to 
    $$a\bmat{-1 \\ 1 \\ 1 \\ 1} + b\bmat{1 \\ -1 \\ 1 \\ 1} + c\bmat{1 \\ 1 \\ -1 \\ 1} + d\bmat{1 \\ 1 \\ 1 \\ -1} = \bmat{1 \\ 8 \\ -7 \\ 2}$$

    The vector $\bmat{1 & 8 & -7 &2}^{\mathsf{T}}$ is a linear combination of the orthogonal set $B$. By the preceding theorem,
    $$c = \frac{\langle\bmat{1 & 8 & -7 & 2}^{\mathsf{T}}, \bmat{1 & 1 & -1 & 1}^{\mathsf{T}}\rangle}{\langle\bmat{1 & 1 & -1 & 1}^{\mathsf{T}}, \bmat{1 & 1 & -1 & 1}^{\mathsf{T}}\rangle} = \frac{18}{4} = \boxed{\frac{9}{2}}$$
\end{enumerate}
\end{example}

\begin{remark}
If $A \in M_n(\RR)$ has orthogonal columns $a_1, a_2, \dots, a_n$, then the matrix equation $Ax = b$ has a unique solution $x = \bmat{x_1 & x_2 \cdots & x_n}^{\mathsf{T}}$ given by
$$x_i = \frac{\langle b, a_i\rangle}{\langle a_i, a_i \rangle}$$
\end{remark}

\begin{theorem}
If $A \in M_n(\RR)$ has orthogonal columns, then $A$ is nonsingular (If $A$ is orthonormal, then $A$ is nonsingular and its inverse $A\inv = A^{\mathsf{T}}$).
\end{theorem}

\begin{definition}
Let $u, v \in V$ with $v \neq \vec{0}$ and $u \neq \operatorname{span}\{v\}$. The vector
$$\operatorname{proj}_vu = \frac{\langle u, v \rangle}{\langle v, v \rangle}v$$
is called the (orthogonal) projection of $u$ onto $v$.
\end{definition}

\begin{remark}
\leavevmode
\begin{enumerate}
    \item $\operatorname{proj}_vu \in \operatorname{span}\{v\}$
    \item $u - \operatorname{proj}_vu \perp v$
\end{enumerate}
\end{remark}

\begin{definition}
Let $S = \{w_1, w_2, \dots, w_n\}$ be an orthogonal subset of $V$ such that $\vec{0} \notin S$ and let $u \notin \operatorname{span}S$. The vector
$$\operatorname{proj}_Su = \sum_{i=1}^n \operatorname{proj}_{w_i}u = \sum_{i=1}^n \frac{\langle u, w_i\rangle}{\langle w_i, w_i \rangle}w_i$$
is called the \textbf{orthogonal projection} of $u$ onto $S$.
\end{definition}

\begin{remark}
\leavevmode
\begin{enumerate}
    \item $\operatorname{proj}_Su \in \operatorname{span}S$
    \item $u - \operatorname{proj}_Su \perp S$
\end{enumerate}
\end{remark}
\newpage
\section{Gram-Schmidt Orthonormalization Process}
\begin{definition}[Gram-Schmidt Orthonormalization]
Let $B = \{v_1, v_2, \dots, v_n\}$ be a basis for $V$.
\begin{enumerate}
    \item Put $w_1 = v_1$.
    \item Put $w_2 = v_2 - \operatorname{proj}_{w_1}v_2$.
    \item Put $w_3 = v_3 - \operatorname{proj}_{\{w_1, w_2\}}v_3$
    \item Combine the above process to get an orthogonal set $$\{w_1, w_2, \dots, w_n\} = S$$
    \item Put $w_n = v_n - \operatorname{proj}_Sv_n$,
    $$B' = \{w_1, w_2, \dots, w_n\}$$
    \item Put $b_i = \dfrac{w_i}{\lVert w_i \rVert}$. Then $B'' = \{b_1, b_2, \dots, b_n\}$ is an orthonormal basis for $V$.
\end{enumerate}
\end{definition}

\begin{remark}
\leavevmode
\begin{enumerate}
    \item The output of a Gram-Schmidt Process is not unique. Hence a permutation $\pi$ could lead to a different basis.
    \item The recursive definition of the process is:
    $$w_k = v_k - \sum_{i = 1}^{k-1}\frac{\langle v_k, w_i\rangle}{\langle w_i, w_i \rangle}w_i$$
\end{enumerate}
\end{remark}

\begin{theorem}
Every finite-dimensional inner product space $V$ has an orthonormal basis.
\end{theorem}

\begin{example}
Transform the basis $S = \{(1, 1, 1), (1, -1, -1), (0, 1, 0)\}$ for $\RR^3$ into an orthonormal basis under the standard inner product.
\end{example}

\textit{Solution.} Let us apply Gram-Schmidt Process,
\begin{itemize}
    \item Put $w_1 = v_1 = (1,1, 1)$
    \item Put $w_2 = v_2 - \operatorname{proj}_{w_1}v_2$ 
    $$= v_2 - \frac{\langle w_1, v_2\rangle}{\langle w_1, w_1\rangle}w_1$$
    $$=(1, -1, -1) + \frac{1}{3}(1, 1, 1) = \left(\frac{4}{3}, -\frac{2}{3}, -\frac{2}{3}\right)$$
    \item Put $w_3 = v_3 - \operatorname{proj}_{w_1}v_3 - \operatorname{proj}_{w_2}v_3$
    $$=v_3 - \frac{\langle w_1, v_3\rangle}{\langle w_1, w_1 \rangle}w_1 - \frac{\langle w_2, v_3 \rangle}{\langle w_2, w_2 \rangle}w_2$$
    $$=(0, 1, 0) - \frac{1}{3}(1, 1, 1) + \frac{1}{4}\left(\frac{4}{3}, -\frac{2}{3}, -\frac{2}{3} \right)$$
    $$=\left(0, \frac{1}{2}, -\frac{1}{2}\right)$$
\end{itemize}

Normalizing each $w_i$,
\begin{itemize}
    \item $w_1 = \dfrac{1}{\sqrt{3}}(1, 1, 1)$
    \item $w_2 = \sqrt{\dfrac{3}{8}}\left(\dfrac{4}{3}, -\dfrac{2}{3}, -\dfrac{2}{3}\right)$
    \item $w_3 = \sqrt{2}\left(0, \dfrac{1}{2}, -\dfrac{1}{2}\right)$
\end{itemize}

Therefore, $\displaystyle B = \left\{\frac{1}{\sqrt{3}}(1, 1, 1), \sqrt{\frac{3}{8}}\left(\frac{4}{3}, -\frac{2}{3}, -\frac{2}{3}\right), \sqrt{2}\left(0, \frac{1}{2}, -\frac{1}{2}\right)\right\}$ is an orthonormal basis for $\RR^3$.

\newpage
\section{Linear Transformations}
\begin{definition}
Let $U$ and $V$, be vector spaces. A \textbf{linear transformation} from $U$ to $V$ is a function $T: U \to V$ satisfying the following:
\begin{enumerate}
    \item $\forall u_1, u_2 \in U, \quad T(u_1 + u_2) = T(u_1) + T(u_2)$ 
    \item $\forall u \in U, \forall c \in \RR, \quad T(cu) = c \cdot T(u)$
\end{enumerate}
\end{definition}

\begin{remark}
\leavevmode
\begin{enumerate}
    \item If $U = V$, then $T$ is called a \textbf{linear operator} on $V$.
    \item Let $T : U \to V$ be a linear transformation
    \begin{enumerate}
        \item $\forall u_1, u_2 \in U, \quad T(u_1 - u_2) = T(u_1) - T(u_2)$
        \item Denote by $\vec{0}_u$ and $\vec{0}_v$ the respective zero vectors of $U$ and $V$. Then $T(\vec{0}_u) = \vec{0}_v$.
    \end{enumerate}
    Therefore, $T(\vec{0}_u + \vec{0}_u) = \vec{0_v} = T(\vec{0}_u)$
    \item $T: U \to V$ is a linear transformation iff $\forall u_1, u_2 \in U, \forall a, b \in \RR$,
    $$T(au_1 + bu_2) = aT(u_1) + bT(u_2)$$
\end{enumerate}
\end{remark}

\begin{example}
Let $U$ and $V$ be vector spaces. The following are linear transformations:
\begin{enumerate}
    \item $T: U \to V$ given by $T(u) = \vec{0}_v, \forall u \in U$.
    \leavevmode
    
    \textit{Solution.} Let $u_1, u_2 \in U$ and $c \in \RR$. Then,
    $$T(u_1 + u_2) = \vec{0}_v$$
    $$T(u_1) + T(u_2) = \vec{0}_v + \vec{0}_v = \vec{0}_v = T(u_1 + u_2)$$
    $$T(cu_1) = \vec{0}_v$$
    $$c \cdot T(u_1) = c\vec{0}_v = \vec{0}_v = T(cu_1)$$
    Therefore, $T$ is a linear transformation.
    \item $\operatorname{id}_u: U \to U$ given by $\operatorname{id}_u(u) = u$
    \leavevmode

    \textit{Solution.} Let $u_1, u_2 \in U$ and $a, b \in \RR$. Then,
    $$\operatorname{id}_u(au_1 + bu_2) = au_1 + bu_2$$
    $$=a\operatorname{id}_u(u_1) + b\operatorname{id}_u(u_2) $$
    Therefore, $\operatorname{id}_u$ is a linear transformation.
\end{enumerate}
\end{example}

\begin{example}
Let $T: \RR^3 \to \RR^2$ be defined by $T(x, y, z) = (x, y)$.

\textit{Solution.} Let $u_1 = (x_1, y_1, z_1), u_2 = (x_2, y_2, z_2) \in \RR^3$ and $a, b \in \RR$. Then,
$$T(au_1 + bu_2) = T(ax_1 + bx_2, ay_1 + by_2, az_1 + bz_2)$$
$$=(ax_1 + bx_2, ay_1 + by_2)$$
$$= a(x_1, y_1) + b(x_2, y_2)$$
$$=aT(x_1, y_1, z_1) + bT(x_2, y_2, z_2)$$
$$= aT(u_1) + bT(u_2)$$
Therefore, $T$ is a linear transformation.
\end{example}

\newpage
\begin{example}
Let $A \in \matmn$ and $T_A: M_{n \times 1}(\RR) \to M_{m \times 1}(\RR)$ be defined by $T_A(x) = Ax$.

\textit{Solution.} Let $x, y \in M_{n \times 1}(\RR)$ and $a, b \in \RR$. Then,
$$T_A(ax + by) = A(ax + by)$$
$$=aAx + bAy$$
$$=aT_A(x) + bT_A(y)$$
Therefore, $T_A$ is a linear transformation.
\end{example}

\begin{example}
Let $n \in \ZZ^+$ and $D: P_n \to P_{n-1}$, be defined by $D(p(x)) = p\prime(x)$

\textit{Solution.} Let $p(x), q(x) \in P_n$ and $c, d \in \RR$. Then
$$D(cp(x) + dq(x)) = \frac{d}{dx}(cp(x) + dq(x))$$
$$=cp'(x) + dq'(x)$$
$$=cD(p(x)) + dD(q(x))$$
Therefore, $D$ is a linear transformation.
\end{example}

\noindent\rule{\textwidth}{0.4pt}

Now, let us consider some examples which uses counterexample to prove that a function is \textbf{NOT} a linear transformation.

\begin{example}
Let $T: \RR \to \RR$ given by $T(x) = |x|$.

\textit{Solution.} Take $u_1 = 2, u_2 = -2$. Then,
$$T(u_1 + u_2) = T(0) = |0| = 0$$
$$T(u_1) + T(u_2) = |-2| + |2| = 4$$
$$\because T(u_1 + u_2) \neq T(u_1) + T(u_2)$$
Therefore, $T$ is not a linear transformation.
\end{example}

\begin{example}
Let $T: \RR^2 \to P_2$ given by $T(a, b) = x^2 + ax + b$.

\textit{Solution.} Let $(a, b) = (6, 7)$ and $c = 2$. Then,
$$T(c(6, 7)) = T(6c, 7c) = T(12, 14) = x^2 + 12x + 14$$
$$cT(6, 7) = 2(x^2 + 6x + 7) = 2x^2 + 12x + 14$$
$$\because T(cu_1) \neq cT(u_1)$$
Therefore, $T$ is not a linear transformation.
\end{example}

\noindent\rule{\textwidth}{0.4pt}

Now suppose we flip it around, and we would like to know $T(u)$ where $T$ is a linear transformation, $U$ is a vector space, and $u \in U$. How can we derive this? It turns out that we can do so by rewriting a vector $u \in U$ as a \textbf{linear combination}! Recall that $u = c_1u_1 + c_2u_2 + \cdots + c_nu_n$ for some $c_1, c_2, \dots, c_n \in \RR$.
$$T(u) = T(c_1u_1 + c_2u_2 + \cdots + c_nu_n)$$
$$=c_1T(u_1) + c_2T(u_2) + \cdots + c_nT(u_n)$$
Therefore, $T(u_1), T(u_2), \dots, T(u_n)$ are sufficient to determine $T(u)$. Then, $T$ is said to be \textbf{completely determined} by $T(u_1), T(u_2), \dots, T(u_n)$.


\newpage
\begin{example}
Let $T: \RR^2 \to \RR^2$ be a linear transformation such that $T(1, 0) = (2, 1)$ and $T(0, 1) = (-3, 2)$. Find $T(x, y)$ for any $x, y \in \RR$.

\textit{Solution.} Let $(x, y) \in \RR^2$. Then,
$$T(x, y) = T(x(1, 0) + y(0, 1))$$
$$=xT(1, 0) + yT(0, 1)$$
$$=x(2, 1) + y(-3, 2)$$
$$=\boxed{(2x - 3y, x + 2y)}$$
\end{example}

\begin{definition}
Let $T: U \to V$ be a linear transformation. 
\begin{enumerate}
    \item The \textbf{kernel} of $T$ is $\ker T = \{u \in U \mid T(u) = \vec{0}_v\} \subseteq U$.
    \item The \textbf{image/range} of $T$ is $\im T = \{T(u) \mid u \in U\} \subseteq V$.
\end{enumerate}
\end{definition}

\begin{remark}
Let $u \in U$ and $v \in V$.
\begin{enumerate}
    \item $u \in \ker T$ iff $T(u) = \vec{0}_v$
    \item $v \in \im T$ iff $v = T(u)$ for some $u \in U$
\end{enumerate}
\end{remark}

\begin{theorem}
Let $T: U \to V$ be a linear transformation. Then $\ker T \leq U$ and $\im T \leq V$.
\end{theorem}

\begin{definition}
Let $T: U \to V$ be a linear transformation with $U$ and $V$ be finite-dimensional.
\begin{enumerate}
    \item The \textbf{rank} of $T$ is $\rank T = \dim(\im T)$
    \item The \textbf{nullity} of $T$ is $\upsilon(T) = \dim(\ker T)$
\end{enumerate}
\end{definition}

\begin{example}
Let $A \in \matmn$ and $T_A: M_{n \times 1}(\RR) \to M_{m \times 1}(\RR)$ defined by $T_A(x) = Ax$.
\begin{enumerate}[label=(\alph*)]
    \item 
    \begin{align*}
    \ker T_A = \{x \in M_{n \times 1}(\RR) \mid T_A(x) = 0_{m \times 1}\} \\
    =\{x \in M_{n \times 1}(\RR) \mid Ax = 0_{m \times 1}\} \\
    =N(A), \quad \text{the null space of $A$} \\
    \therefore \upsilon(T_A) = \dim(\ker T_A) = \dim N(A) = \upsilon(A)
    \end{align*}
    
    
    \item 
    $$\im T_A = \{T_A(x) \mid x \in M_{n \times 1}(\RR)\}$$
    $$=\{Ax \mid x \in M_{n \times 1}(\RR)\}$$
    $$=\{x_1C_1 + x_2C_2 + \cdots + x_nC_n \mid x_1, x_2, \dots, x_n \in \RR \hspace{0.2cm} \text{and} \hspace{0.2cm} C_i = \text{col $i$ of $A$}\}$$
    $$=\operatorname{span}\{C_1, C_2, \dots, C_n\}$$
    $$= C(A), \quad \text{the column space of $A$}$$
    $$\therefore \rank(T_A) = \dim(\im T) = \dim C(A) = \rank A$$
\end{enumerate}
\end{example}

\newpage
\begin{example}
Let $T: \RR^3 \to \RR^3$ defined by $T(x, y, z) = (x + z, x + y, y - z)$.\footnote{As an exercise prove that this is a linear transformation.} Find the bases for the image and kernel of $T$.

\textit{Solution.}

\begin{enumerate}[label=(\alph*)]
    \item Recall that $(x, y, z) \in \ker T$ iff $T(x, y, z) = (0, 0, 0)$. Then,
    $$(x + z, x + y, y - z) = (0, 0, 0) \quad \longleftrightarrow \quad \begin{cases}
        x + z = 0 \\
        x + y = 0 \\
        y - z = 0
    \end{cases}$$
    Turning this into a coefficient matrix,
    $$\bmat{1 & 0 & 1 \\ 1 & 1 & 0 \\ 0 & 1 & -1} \xrightarrow{\text{RREF}} \bmat{1 & 0 & 1 \\ 0 & 1 & -1 \\ 0 & 0 & 0} \longleftrightarrow \begin{cases}
        x + z = 0 \\
        y - z = 0
    \end{cases}$$
    Since the third column is a free variable, we let $z = t$. Then, $x = -t$ and $y = t$. Therefore $\ker T = \{(-t, t, t) \mid t \in \RR\} = \spn\{(-1, 1, 1)\}$. 
    
    Then $B = \{(-1, 1, 1)\}$ is a basis for $\ker T$, and $\upsilon(T) = 1$.

    \item Let us find a basis for $\im T$,
    $$\im T = \{T(x, y, z) \mid x, y, z \in \RR\}$$
    $$=\{(x + z, x + y, y - z) \mid x, y, z \in \RR\}$$
    $$=\{x(1, 1, 0) + y(0, 1, 1) + z(1, 0, -1) \mid x, y, z \in \RR\}$$
    $$=\spn\{(1, 1, 0), (0, 1, 1), (1, 0, -1)\}$$
    Then, let us check whether this is linearly independent. Form,
    $$x(1, 1, 0) + y(0, 1, 1) + z(1, 0, -1) = (0, 0, 0)$$
    The coefficient matrix is,
    $$\bmat{1 & 0 & 1 \\ 1 & 1 & 0 \\ 0 & 1 & -1} \xrightarrow{\text{RREF}} \bmat{1 & 0 & 1 \\ 0 & 1 & -1 \\ 0 & 0 & 0}$$
    Since there is a free variable in the third column, we must delete that. 
    
    Therefore $B = \{(1, 1, 0), (0, 1, 1)\}$ is a basis for $\im T$ and $\rank T = 2$
\end{enumerate}
\end{example}

\leavevmode

\begin{theorem}[Rank-Nullity Theorem]
Let $U$ and $V$ be vector spaces with $\dim U = n < \infty$. If $T: U \to V$ is a linear transformation, then
$$\dim U = \rank T + \upsilon(T)$$
\end{theorem}

\begin{definition}
Let $T: U \to V$ be a linear transformation.
\begin{enumerate}
    \item $T$ is \textbf{one-to-one} if $\forall u_1, u_2 \in U$. $T(u_1) = T(u_2)$ iff $u_1 = u_2$.
    \item $T$ is \textbf{onto} if $\im T = V$, i.e. $\forall v \in V, \hspace{0.2cm} \exists u \in U$ such that $v = T(u)$.
    \item $T$ is a \textbf{bijection/isomorphism} if $T$ is both one-to-one and onto.
\end{enumerate}
\end{definition}

\newpage
\begin{remark}
Let $U$ and $V$ be vector spaces,
\begin{enumerate}
    \item If there exists an isomorphism $T: U \to V$, then $U$ and $V$ are said to be isomorphic, denoted by $U \cong V$.
    \item If $T: U \to V$ is a linear transformation, then $T$ is one-to-one iff $\ker T = \{\vec{0}_v\}$. Only contains the zero vector.
    \item $T$ is onto iff $\rank T = \dim V$.   
\end{enumerate}
\end{remark}

\lline

For the next few examples, we will \textbf{determine the mapping} of a given function using the remarks given above.

\begin{example}
Determine the mapping of a function $\id_U: U \to U$ defined by $\id_U(u) = u$.
    
\textit{Solution.} 

\begin{enumerate}
    \item $u \in \ker(\id_U)$ iff $\id_U(u) = \vec{0}_u$ and $u = \vec{0}_u$. 
    
    Therefore, $\ker(\id_U) = \vec{0}_u$, hence, $\id_U$ is one-to-one.

    \item $\im(\id_U) = \{\id_U(u) \mid u \in U\} = \{u \mid u \in U\}$

    Therefore $\id_U$ is onto. And since $\id_U$ is both one-to-one and onto, it is an isomorphism.
\end{enumerate}
\end{example}

\begin{example}
Determine the mapping of a function $T: \RR^2 \to \RR^3$ defined by $T(x, y) = (x, y, 0)$.

\textit{Solution.}

\begin{enumerate}[label=(\alph*)]
    \item $(x, y) \in \ker T$ iff $T(x, y) = (0, 0, 0)$
    $$(x, y, 0) = (0, 0, 0) \Longrightarrow x = y = 0$$
    Therefore, $\ker T = \{\vec{0}\}$, hence, $T$ is one-to-one.

    \item $\im T = \{T(x, y) \mid x, y \in \RR\}$
    $$=\{(x, y, 0) \mid x, y \in \RR\}$$
    $$=\{x(1, 0, 0) + y(0, 1, 0) \mid x, y \in \RR\}$$
    $$=\spn\{(1, 0, 0), (0, 1, 0)\}$$
    Since neither of them are scalar multiples of each other it is a basis for $\im T$. And since $\rank T \neq \dim \RR^3$, it is not onto.

    Therefore, $T$ is one-to-one.
\end{enumerate}
\end{example}

\begin{example}
Determine the mapping of a function $T:\RR^3 \to P_2$ defined by $T(a, b, c) = (a + b)x^2 + (b + c)x + c$.

\textit{Solution.}

\begin{enumerate}[label=(\alph*)]
    \item $(a, b, c) \in \ker T$ iff $(a + b)x^2 + (b + c)x + c = 0$
    $$\begin{cases}
        a + b = 0 \\
        b + c = 0 \\
        c = 0
    \end{cases} \quad \longleftrightarrow \quad (a, b, c) = (0, 0, 0)$$
    Therefore $\ker T = \{(0, 0, 0)\}$, hence, $T$ is one-to-one.

    \item $\im T = \{T(a, b, c) \mid a, b, c \in \RR\}$
    $$=\{(a + b)x^2 + (b + c)x + c \mid a, b, c \in \RR\}$$
    $$=\{a(x^2) + b(x^2 + x) + c(x + 1) \mid a, b, c \in \RR\}$$
    $$=\spn\{\underbrace{x^2, x^2 + x, x + 1}_B\}$$
    We can verify that $B$ is a basis for $\im T$. And since $\rank T = 3 = \dim P_2$, hence, $T$ is onto.

    Therefore, $T$ is an isomorphism, and $\RR^3 \cong P_2$.
\end{enumerate}
\end{example}

\begin{theorem}
Let $U, V, W$ be vector spaces and $T: U \to V$ and $S: V \to W$ be linear transformations.
\begin{enumerate}
    \item The function $S \circ T: U \to W$ is a linear transformation. If $S$ and $T$ are isomorphic, then so is $S \circ T$.
    \item $T$ is an isomorphism iff there exists a unique function $T\inv: V \to U$ such that,
    $$(T \circ T\inv)(v) = v, \quad \forall v \in V$$
    $$(T\inv \circ T)(u) = u, \quad \forall u \in U$$
\end{enumerate}
\end{theorem}

\begin{remark}
Let $U, V, W$ be vector spaces
\begin{enumerate}
    \item If $T: U \to V$ is an isomorphism, then the function $T\inv: V \to U$ is also an isomorphism, called the inverse of $T$.
    \item
    \begin{enumerate}
        \item $U \cong U$
        \item If $U \cong V$, then $V \cong U$
        \item If $U \cong V$ and $V \cong W$, then $U \cong W$
    \end{enumerate}
\end{enumerate}
\end{remark}

\begin{theorem}
Let $U$ and $V$ be finite-dimensional vector spaces.
\begin{enumerate}
    \item If $\dim V = n$, then $V \cong \RR^n$
    \item $U \cong V$ iff $\dim U = \dim V$
    \item Let $T: U \to V$ be a linear transformation. If $\dim U = \dim V$, then the following are equivalent:
    \begin{enumerate}
        \item $T$ is one-to-one
        \item $T$ is onto
        \item $T$ is an isomorphism
    \end{enumerate}
\end{enumerate}
\end{theorem}

\begin{example}
\leavevmode
\begin{enumerate}
    \item $M_{n \times 1}(\RR) \cong \RR^n$
    \item $\matmn \cong M_{n \times m}(\RR) \cong \RR^{mn}$
    \item $P_n \cong \RR^{n + 1}$
    \item Show that $T: \RR^2 \to \CC$ given by $T(a, b) = (a + b) + (a - b)i$ is an isomorphism.
\end{enumerate}
\end{example}

\begin{remark}
Let $T: U \to V$ be a linear transformation.
\begin{enumerate}
    \item $T$ is not necessarily and isomorphism even when $U \cong V$. It only guarantees that \ul{such functions exist}.
    \item If $\dim U \neq \dim V$, then $T$ is \textbf{NOT} an isomorphism.
    \begin{enumerate}
        \item If $\dim U > \dim V$, then $T$ is not one-to-one.
        \item If $\dim U < \dim V$, then $T$ is not onto.
    \end{enumerate}
\end{enumerate}
\end{remark}
\newpage
\section{Matrix of Linear Transformations}
\begin{definition}
Let $V$ be a finite-dimensional vector space. An \textbf{ordered basis} for $V$ such that the order in which the elements are listed is taken into account.
\end{definition}

\begin{example}
In $\RR^3$, $B_1 = \{e_1, e_2, e_3\}$ and $B_2 = \{e_3, e_1, e_2\}$ are equal sets, but not ordered bases.
\end{example}

\begin{example}
Let $V = \RR^2$ with ordered bases
$$B_1 = \{(1, 0), (0, 1)\}$$
$$B_2 = \{(0, 1), (1, 0)\}$$
$$B_3 = \{(1, 2), (3, 5)\}$$
\begin{enumerate}
    \item If $v = (-2, 5)$, find $[v]_{B_i}$, $i \in \{1, 2, 3\}$

    \textit{Solution.} 
    \leavevmode
    \begin{itemize}
        \item $v = (-2, 5) = -2(1, 0) + 5(0, 1)$
        $$[v]_{B_1} = \bmat{-2 \\ 5}$$
        $$[v]_{B_2} = \bmat{5 \\ -2}$$

        \item To find $[v]_{B_3}$,
        $$v = (-2, 5) = x(1, 2) + y(3, 5) = (x + 3y, 2x + 5y)$$
        $$\begin{amatrix}{2}
            1 & 3 & -2 \\
            2 & 5 & 5 \\
        \end{amatrix} \xrightarrow{\text{RREF}} 
        \begin{amatrix}{2}
            1 & 0 & 25 \\
            0 & 1 & -9 \\
        \end{amatrix}$$
        $$[v]_{B_3} = \bmat{25 \\ -9}$$
    \end{itemize}

    \item If $w \in \RR^2$ such that $[w]_{B_3} = \bmat{-1 \\ 3}$, find $w$.

    \textit{Solution.} 
    $$w = -1(1, 2) + 3(3, 5) = \boxed{(8, 13)}$$
\end{enumerate}
\end{example}

\begin{definition}
Let $U$ and $V$ be vector spaces with respective ordered bases
$$B_U = \{u_1, u_2, \dots, u_n\}$$
$$B_V = \{v_1, v_2, \dots, v_m\}$$
If $T: U \to V$ is a linear transformation, the \textbf{matrix representation of $T$ with respect to $B_U$ and $B_V$} is
$$[T]_{B_U}^{B_V} = \bmat{[T(u_1)]_{B_V} & [T(u_2)]_{B_V} & \cdots & [T(u_n)]_{B_V}} \in M_{m \times n}(\RR)$$
\end{definition}

\begin{example}
Let $V$ be a $3$-dimensional vector space with ordered basis $B = \{v_1, v_2, v_3\}$. If $\id_V: V \to V$ is the identity linear transformation, find $[\id_V]_B^B$.

\textit{Solution.} 
\begin{align*}
[\id_V]_B^B &= \bmat{[\id_V(v_1)]_B & [\id_V(v_2)]_B & [\id_V(v_3)]_B} \\
&=\bmat{[v_1]_B & [v_2]_B & [v_3]_B}
\end{align*}
\end{example}

\newpage
Then we must find the values in $B$ under $\id_V$,
\begin{align*}
v_1 &= 1v_1 + 0v_2 + 0v_3 \\
v_2 &= 0v_1 + 1v_2 + 0v_3 \\
v_3 &= 0v_1 + 0v_2 + 1v_3
\end{align*}
$$\therefore [\id_V]_B^B = \bmat{1 & 0 & 0 \\ 0 & 1 & 0 \\ 0 & 0 & 1}$$

\begin{example}
Let $B_2 = \{(1, 3), (-2, 4)\}$ and $B_3 = \{(1, 1, 1), (2, 2, 0),(3,0,0)\}$ be the ordered bases for $\RR^2$ and $\RR^3$ respectively. If $T: \RR^2 \to \RR^3$ is the linear transformation given by
$$T(x, y) = (x + 2y, -x, 0)$$
Find $[T]_{B_2}^{B_3}$.
\end{example}

\begin{remark}
Let $U$ and $V$ be finite-dimensional vector spaces with respective ordered bases $B_U = \{u_1, u_2, \dots, u_n\}$ and $B_V = \{v_1, v_2, \dots, v_m\}$ and $T: U \to V$ be e linear transformation. To find $[T]_{B_U}^{B_V}$:
\begin{enumerate}
    \item Find $T(u_1), T(u_2), \dots, T(u_n)$
    \item Form the augmented matrix
    $$X = [B_V \mid T(B_U)] = \bmat{v_1 & v_2 & \cdots & v_n & \vline & T(u_1) & T(u_2) & \cdots & T(u_n)}$$
    \item Use elementary row operations to transform $X$ into the matrix $[I_n \mid M]$. Then $M = [T]_{B_U}^{B_V}$.
\end{enumerate}
\end{remark}

\begin{theorem}
Let $U$ and $V$ be finite-dimensional vector spaces with ordered bases $B_U$ and $B_V$, respectively. If $T: U \to V$ is a linear transformation and $u \in U$, then
$$[T(u)]_{B_V} = [T]_{B_U}^{B_V}[u]_{B_U}$$
\end{theorem}

\begin{example}
Consider the ordered bases $B = \{(1, 0, 0), (0, 1, 0), (0, 0, 1)\}$ and $C = \{x^2, x, 1\}$ for $\mathbb{R}^3$ and $P_2$, respectively. If $T : \mathbb{R}^3 \to P_2$ is the linear transformation such that
$$[T]_B^C = \bmat{0 & 1 & 1 \\ 1 & 0 & 1 \\ 1 & 1 & 0}$$
determine $T(a, b, c)$ for all $(a, b, c) \in \RR^3$.
\end{example}

\textit{Solution.} Let $u = (a, b, c) \in \RR^3$. Then
$$[T(a, b, c)]_C = [T]_B^C[(a, b, c)]_B = \bmat{0 & 1 & 1 \\ 1 & 0 & 1 \\ 1 & 1 & 0}\bmat{a \\ b \\ c} = \bmat{b + c \\ a + c \\ a + b}$$
Therefore, $T(a, b, c) = (b + c)x^2 + (a + c)x + (a + b)$ for all $a,b,c \in \RR$.

\begin{theorem}
Let $U, V$, and $W$ be finite-dimensional vector spaces with respective ordered bases $B_U, B_V$, and $B_W$. If $T: U \to V$ and $S: V \to W$ are linear transformations, then
\begin{enumerate}
    \item $[S \circ T]_{B_U}^{B_W} = [S]_{B_V}^{B_W}[T]_{B_U}^{B_V}$
\end{enumerate}
\end{theorem}
\newpage
\section{Eigenvalues and Eigenvectors}

\begin{definition}
Let $A \in M_n(\RR)$. Then $\lambda \in \RR$ is called an \textbf{eigenvalue} of $A$ if there exists nonzero $v \in M_{n \times 1}(\RR)$ such that
$$Av = \lambda v$$
Such a vector $v$ is called an \textbf{eigenvector} of $A$ associated to $A$.
\end{definition}

\begin{remark}
Let $A \in M_n(\RR)$
\begin{enumerate}
    \item We impose the condition that $v \notin 0_{n \times 1}$ since $A0_{n \times 1} = \lambda 0_{n \times 1}$ is true for any $\lambda \in \RR$.
    \item $v$ is an eigenvector of $A$ iff $v$ is a scalar multiple of $Av$.
\end{enumerate}
\end{remark}

\begin{example}
\leavevmode
\begin{enumerate}
    \item Let $v \in M_{n \times 1}(\RR)$ be nonzero 
    \begin{itemize}
        \item $I_nv = 1v$ \\
        Therefore, $\lambda = 1$ is an eigenvalue of $I_n$ where associated eigenvectors are any nonzero $v \in M_{n \times 1}(\RR)$
        \item $0_{n \times n}v = 0_{n \times 1} = 0v$ \\
        Therefore, $\lambda = 0$ is an eigenvalue of $0_{n \times n}$ where associated eigenvectors are nonzero $v \in M_{n \times 1}(\RR)$.
    \end{itemize}
    \item Let $A = \bmat{0 & 2 \\ 2 & 0}$
    \begin{itemize}
        \item $A\bmat{1 \\ 1} = \bmat{2 \\ 2} = 2\bmat{1 \\ 1}$ \\
        $\therefore \lambda_1 = 2$ is an eigenvalue of $A$ with associated eigenvector $v_1 =\bmat{1 \\ 1}$.
        \item $A\bmat{-1 \\ 1} = \bmat{2 \\ -2} = -2\bmat{1 \\ -1}$ \\
        $\therefore \lambda_2 = -2$ is an eigenvalue of $A$ with associated eigenvector $v_2 = \bmat{-1 \\ 1}$.
    \end{itemize}
    \item Let $A = \bmat{0 & 2 & 2 \\ 2 & 0 & 2 \\ 2 & 2 & 0}$.
    \begin{itemize}
        \item $A\bmat{1 \\ 1 \\ 1} = \bmat{4 \\ 4 \\ 4} = 4\bmat{1 \\ 1 \\ 1}$ \\
        $\therefore \lambda_1 = 4$ is an eigenvalue of $A$ with associated eigenvector $v_1 = \bmat{1 \\ 1 \\ 1}$.
        \item $A\bmat{1 \\ 0 \\ -1} = \bmat{-2 \\ 0 \\ 2} = -2\bmat{1 \\ 0 \\ -1}$ \\
        $\therefore \lambda_2 = -2$ is an eigenvalue of $A$ with associated eigenvector $v_2 = \bmat{1 \\ 0 \\ -2}$.
    \end{itemize}
\end{enumerate}
\end{example}

\newpage
\begin{definition}
Let $A \in M_n(\RR)$.
\begin{enumerate}
    \item The \textbf{characteristic polynomial} of $A$ is 
    $$p_A(x) = \det(A - xI_n)$$
    $$=\vmat{a_{11} - x & a_{12} & \cdots & a_{1n} \\ a_{21} & a_{22} - x & \cdots & a_{2n} \\ \vdots & \vdots & \ddots & \vdots \\ a_{n1} & a_{n2} & \cdots & a_{nn} - x}$$
    \item The \textbf{characteristic equation} of $A$ is $p_A(x) = 0$.
\end{enumerate}
\end{definition}

\begin{remark}
Let $A \in M_n(\RR)$.
\leavevmode
\begin{enumerate}
    \item $p_A(x)$ is a polynomial of degree $n$ with leading coefficients $(-1)^n$ and constant term $\det A$.
    \item $\lambda \in \RR$ is an eigenvalue of $A$ iff $p_A(\lambda) = 0$
    \item If $p_A(x)$ has no real roots, then $A$ has no (real) eigenvalues. 
    \item The \textbf{algebraic multiplicity} of $\lambda$ is the multiplicity of $\lambda$ as a root of $p_A(x)$, denoted by $a(\lambda)$
\end{enumerate}
\end{remark}

\begin{definition}
Let $A \in M_n(\RR)$ and $\lambda \in \RR$ be an eigenvalue of $A$. Then
$$E(A, \lambda):= N(A - \lambda I_n)$$
is called the \textbf{eigenspace of $A$ associated to $\lambda$}.
\end{definition}

\begin{remark}
\leavevmode
\begin{enumerate}
    \item $E(A, \lambda)$ is an eigenspace of $M_{n \times 1}(\RR)$. The \textbf{geometric multiplicity} of $\lambda$ is $g(\lambda) = \dim E(A, \lambda) = \upsilon(A - \lambda I_n)$.
    \item The eigenvectors of $A$ associated to $\lambda$ are the nonzero vectors in $E(A, \lambda)$. Thus, $\mu \in \RR$ is NOT an eigenvalue of $A$ iff $N(A - \mu I_n) = \{\vec{0}\}$.
    \item To find (a basis for) $E(A, \lambda)$, solve for $v$ in $(A - \lambda I_n)_v = 0_{n \times 1}$ (by GJR/GE).
\end{enumerate}
\end{remark}

\begin{example}
Find the eigenvalues and corresponding eigenspaces of a matrix $A$ where,
$$A = \bmat{0 & 2 \\ 2 & 0}$$
and identify the algebraic and geometric multiplicity of each eigenvalue.
\end{example}

\textit{Solution.}

Let us first solve for its characteristic polynomial $p_A(x)$,
$$p_A(x) = \vmat{-x & 2 \\ 2 & -x} = x^2 - 4$$
Therefore, the eigenvalues of $A$: $2, -2$

\begin{itemize}
    \item $E(A, 2)$ \\
    Forming the augmented matrix,
    $$\begin{amatrix}{2}
        -2 & 2 & 0 \\
        2 & -2 & 0 \\
    \end{amatrix} \xrightarrow{\text{RREF}}
    \begin{amatrix}{2}
        1 & -1 & 0 \\
        0 & 0 & 0 \\
    \end{amatrix} \longleftrightarrow
    a - b = 0
    $$
    Since there are no leading entries in the second column, let $b = t$, then, $a = t$,
    $$E(A, 2) = \left\{\bmat{t \\ t} \suchthat t \in \RR\right\} = \spn\left\{\bmat{1 \\ 1}\right\}$$
    Since the set only contains one element, it is a basis for $E(A, 2)$. And $a(2) = 1 = g(2)$.

    \item $E(A, -2)$ \\
    Forming the augmented matrix,
    $$\begin{amatrix}{2}
        2 & 2 & 0 \\
        2 & 2 & 0 \\
    \end{amatrix} \xrightarrow{\text{RREF}}
    \begin{amatrix}{2}
        1 & 1 & 0 \\
        0 & 0 & 0 \\
    \end{amatrix} \longleftrightarrow a + b = 0
    $$
    Let $b = t$, then, $a = -t$,
    $$E(A, -2) = \left\{\bmat{-t \\ t} \suchthat t \in \RR\right\} = \spn\left\{\bmat{-1 \\ 1}\right\}$$
    Since the set only contains one element, it is a basis for $E(A, -2)$ and $a(-2) = 1 = g(-2)$.
\end{itemize}

\begin{example}
Find the eigenvalues and corresponding eigenspaces of a matrix $B$ where, 
$$B = \bmat{4 & 2 & 3 \\ 2 & 1 & 2 \\ -1 & -2 & 0}$$
and identify the algebraic and geometric multiplicity of each eigenvalue.
\end{example}

\textit{Solution.} 

Let us first solve for its characteristic polynomial $p_B(x)$,
$$p_B(x) = \vmat{4 - x & 2 & 3 \\ 2 & 1 - x & 2 \\ -1 & -2 & -x} \xrightarrow{\substack{R_1 \leftarrow R_3 + R_1 \\ C_3 \leftarrow C_3 - C_1}} \vmat{3 - x & 0 & 0 \\ 2 & 1 - x & 0 \\ -1 & -2 & 1-x} = (1 - x)^2(3 - x)$$
Therefore, the eigenvalues of $B$: $1, 3$
\begin{itemize}
    \item $E(A, 1)$ \\
    Forming the augmented matrix,
    $$\begin{amatrix}{3}
        3 & 2 & 3 & 0 \\
        2 & 0 & 2 & 0 \\
        -1 & -2 & -1 & 0
    \end{amatrix} \xrightarrow{\text{RREF}}
    \begin{amatrix}{3}
        1 & 0 & 1 & 0 \\
        0 & 1 & 0 & 0 \\
        0 & 0 & 0 & 0
    \end{amatrix} \longleftrightarrow
    \begin{cases}
        a + c = 0 \\ b = 0
    \end{cases}
    $$
    Let $c = t$, then, $a = -t$,
    $$E(A, 1) = \left\{\bmat{-t \\ 0 \\ t} \suchthat t \in \RR\right\} = \spn\left\{\bmat{-1 \\ 0 \\ 1}\right\}$$
    Since it only contains one element it is a basis for the eigenspace. And $a(1) = 2 \neq 1 = g(1)$

    \item $E(A, 3)$ \\
    Forming the augmented matrix,
    $$\begin{amatrix}{3}
        1 & 2 & 3 & 0 \\
        2 & -2 & 2 & 0 \\
        -1 & -2 & -3 & 0 
    \end{amatrix} \xrightarrow{\text{RREF}} 
    \begin{amatrix}{3}
        1 & 2 & 3 & 0 \\
        0 & 1 & \frac{2}{3} & 0 \\
        0 & 0 & 0 & 0
    \end{amatrix} \longleftrightarrow
    \begin{cases}
        a + 2b + 3c = 0 \\
        b + \frac{2}{3}c = 0
    \end{cases}$$
    Let $c = t$, then, $b = -\dfrac{2}{3}t, a = -\dfrac{5}{3}t$,
    $$E(A, 3) = \left\{\bmat{-\frac{5}{3}t \\ -\frac{2}{3}t \\ t} \suchthat t \in \RR\right\} = \spn\left\{\bmat{-\frac{5}{3} \\ -\frac{2}{3} \\ 1} \suchthat t \in \RR\right\}$$
    Since it only contains one element it is a basis, and $a(3) = 1 = g(3)$.
\end{itemize}

\begin{example}
Find the eigenvalues and corresponding eigenspaces of a matrix $C$ where, 
$$C = \bmat{1 & 2 & 0 & 0 \\ 0 & 1 & 0 & 0 \\ 0 & 0 & -1 & 0 \\ 0 & 0 & 5 & 1}$$
and identify the algebraic and geometric multiplicity of each eigenvalue.
\end{example}

Let us first determine the characteristic polynomial $p_C(x)$,
$$p_C(x) = \vmat{1 - x & 2 & 0 & 0 \\ 0 & 1 - x & 0 & 0 \\ 0 & 0 & -1 - x & 0 \\ 0 & 0 & 5 & 1 - x}$$
Performing cofactor expansion along column 1,
$$=(1-x)\vmat{1 - x & 0 & 0 \\ 0 & -1 - x & 0 \\ 0 & 5 & 1- x} = (1-x)^3(-1-x)$$
Therefore, the eigenvalues of $C$: $1, -1$
\begin{itemize}
    \item $E(A, 1)$ \\
    $$\begin{amatrix}{4}
        0 & 2 & 0 & 0 & 0 \\
        0 & 0 & 0 & 0 & 0 \\
        0 & 0 & -2 & 0 & 0 \\
        0 & 0 & 5 & 0 & 0
    \end{amatrix} \xrightarrow{\text{RREF}}
    \begin{amatrix}{4}
        0 & 1 & 0 & 0 & 0 \\ 
        0 & 0 & 1 & 0 & 0 \\
        0 & 0 & 0 & 0 & 0 \\
        0 & 0 & 0 & 0 & 0
    \end{amatrix} \longleftrightarrow
    \begin{cases}
        b = 0 \\
        c = 0
    \end{cases}
    $$
    Let $a = s$ and $d = t$, then,
    $$E(A, 1) = \left\{\bmat{s \\ 0 \\ 0 \\ t} \suchthat s, t \in \RR\right\} = \left\{s\bmat{1 \\ 0 \\ 0 \\ 0} + t \bmat{0 \\ 0 \\ 0 \\ 1} \suchthat s, t \in \RR\right\}$$
    $$=\spn\left\{\bmat{1 \\ 0 \\ 0 \\ 0}, \bmat{0 \\ 0 \\ 0 \\ 1}\right\}$$
    Since they are not scalar multiples of each other, it is a basis for the eigenspace. And $a(1) = 3\neq 2 = g(1)$.
    \item $E(A, -1)$
    $$\begin{amatrix}{4}
        2 & 2 & 0 & 0 & 0 \\
        0 & 2 & 0 & 0 & 0 \\
        0 & 0 & 0 & 0 & 0 \\
        0 & 0 & 5 & 2 & 0 
    \end{amatrix} \xrightarrow{\text{RREF}}
    \begin{amatrix}{4}
        1 & 1 & 0 & 0 & 0 \\
        0 & 1 & 0 & 0 & 0 \\
        0 & 0 & 1 & \frac{2}{5} & 0 \\
        0 & 0 & 0 & 0 & 0
    \end{amatrix} \longleftrightarrow
    \begin{cases}
        a + b = 0 \\
        b = 0 \\
        c + \frac{2}{5}d = 0
    \end{cases}
    $$
    Let $d = t$, then, $a = 0, c = -\dfrac{2}{5}t$,
    $$E(A, -1) = \left\{\bmat{0 \\ 0 \\ -\frac{2}{5}t \\ t}\suchthat t \in \RR\right\} = \spn\left\{\bmat{0 \\ 0 \\ -\frac{2}{5} \\ 1}\right\}$$
    Since it only contains one element, it is a basis for the eigenspace. And $a(-1) = 1 = g(-1)$.
\end{itemize}

\ul{\textbf{Question}.}

Now, given a square matrix when is it nonsingular if we only know its eigenvalue? Let us observe. Suppose that there exists a matrix $A \in M_n(\RR)$. Then, $A$ is nonsingular iff $\det A \neq 0$. This implies that $p_A(0) \neq 0$ and $0$ is not a root of the characteristic polynomial, hence $0$ is not an eigenvalue of $A$.

\begin{theorem}
Let $A \in M_n(\RR)$. Then $A$ is nonsingular iff $0$ is not an eigenvalue of $A$.
\end{theorem}


\newpage
\section{Similarity, Diagonalization, and Symmetric Matrices}
\begin{definition}
Let $A, B \in M_n(\RR)$. Then $A$ and $B$ are \textbf{similar} if there exists $P \in M_n(\RR)$ such that $B = P\inv AP$.
\end{definition}

\begin{remark}
Let $A, B \in M_n(\RR)$ be similar, and $V$ be an $n$-dimensional vector space.
\begin{enumerate}
    \item Then $A$ and $B$ represent the same linear operator $T: U \to V$.
    \item Then $A$ and $B$ are equivalent, hence $\rank A = \rank B$.
\end{enumerate}
\end{remark}

\lline

Now, let's explore some of the properties of similar matrices. Suppose that there exists $A, B \in M_n(\RR)$ such that $A$ and $B$ are similar. Then, it follows that $B = P\inv AP$ for some $P \in M_n(\RR)$.

Since $A$ and $B$ are square matrices, we could obtain their determinant,
\begin{align*}
\det B & = \det(P\inv AP) \\
&=\det(P\inv)\det A \det P \\
&= \frac{1}{\det P}(\det A)(\det P) \\
&=\det A
\end{align*}

Now let's observe their characteristic polynomials,
\begin{align*}
p_B(x) & = \det(B - xI_n) = \det(P\inv AP - xP\inv P) \\
& = \det[(P\inv A - xP\inv(xI_n))P] \\
& = \det[P\inv(A - xI_n)P] \\
& = \det(A - xI_n) \\
& = p_A(x)
\end{align*}


\begin{theorem}
Let $A, B \in M_n(\RR)$ be similar. Then
\begin{enumerate}
    \item $\rank A = \rank B$ and $\upsilon(A) = \upsilon(B)$
    \item $\det A = \det B$
    \item $p_A(x) = p_B(x)$ (hence $A$ and $B$ have the same eigenvalues)
    \item $A^k$ is similar to $B^k$ for all $k \in \ZZ^+$
    \item $A^{\mathsf{T}}$ is similar to $B^{\mathsf{T}}$ 
    \item If $A$ is nonsingular, then $B$ is nonsingular and $A\inv$ is similar to $B\inv$
\end{enumerate}
\end{theorem}

\begin{definition}
$A \in M_n(\RR)$ is said to be \textbf{diagonalizable} of $A$ is similar to a diagonal matrix, i.e. there exists $P \in M_n(\RR)$ such that $P\inv AP$ is diagonal.
\end{definition}

An example of this is suppose we have two matrices $A, P \in M_n(\RR)$, where,
$$A = \bmat{0 & 1 \\ 1 & 0}, \quad P = \bmat{1 & -1 \\ 1 & 1}$$
Then,
$$P\inv AP = \bmat{1 & 0 \\ 0 & -1}$$
Therefore, $A$ is diagonalizable.

\ul{\textbf{Question.}} When is $A \in M_n(\RR)$ diagonalizable?

$A$ is diagonalizable iff there exists a nonsingular matrix $P \in M_n(\RR)$ such that $P\inv AP = D = \diag(d_1, \dots, d_n)$. Let $P_i$ be column $i$ of $P$. Then $AP = PD$. Now


\begin{theorem}
$A \in M_n(\RR)$ is diagonalizable iff $A$ has $n$ linearly independent eigenvectors.
\end{theorem}

\lline

To diagonalize $A \in M_n(\RR)$:
\begin{enumerate}
    \item Find the eigenvalues $\lambda_i$ of $A$.
    \item Find bases for each eigenspace $E(A, \lambda_i)$.
    \item Form the union $E$ of the bases found in (2). Then $E$ is linearly independent. If $|E| = n$, then $A$ is diagonalizable. Otherwise $A$ is not diagonalizable.
    \item If $A$ is diagonalizable, then $P\inv AP = \diag(\lambda_1, \lambda_2, \dots, \lambda_n)$ where $\lambda_1, \lambda_2, \dots, \lambda_n \in \RR$ are the eigenvalues (counting algebraic multiplicities and column $i$ of $P$ is an eigenvector of $A$ associated to $\lambda_i$). We will denote the set of eigenvalues as $\sigma(A)$.
\end{enumerate}

\begin{example}
Determine if the matrix $A$ is diagonalizable. If so, find diagonal $D$ and nonsingular $P$ such that $D = P\inv AP$. Where,
$$A = \bmat{1 & 0 & 0 & 0 \\ 0 & 1 & 5 & -10 \\ 1 & 0 & 2 & 0 \\ 1 & 0 & 0 & 3}$$
\end{example}

\textit{Solution.} 

We first compute for the characteristic polynomial $p_A(x)$,
$$p_A(x) = \vmat{1 - x & 0 & 0 & 0 \\
0 & 1 - x & 5 & -10 \\ 1 & 0 & 2 - x & 0 \\ 1 & 0 & 0 & 3 - x}$$
By performing cofactor expansion along row 1,
$$=1 - x \vmat{1 - x & 5 & -10 \\ 0 & 2 - x & 0 \\ 0 & 0 & 3 - x} =(1-x)^2(2-x)(3-x)$$
Therefore, the eigenvalues of $A$: $1, 2, 3$

\begin{itemize}
    \item $E(A, 1)$
    $$\begin{amatrix}{4}
        0 & 0 & 0 & 0 & 0 \\
        0 & 0 & 5 & -10 & 0 \\
        1 & 0 & 1 & 0 & 0 \\
        1 & 0 & 0 & 2 & 0 
    \end{amatrix} \xrightarrow{\text{RREF}}
    \begin{amatrix}{4}
        1 & 0 & 0 & 2 & 0 \\
        0 & 0 & 1 & -2 & 0 \\
        0 & 0 & 0 & 0 & 0 \\
        0 & 0 & 0 & 0 & 0 
    \end{amatrix} \longleftrightarrow 
    \begin{cases}
        a + 2d = 0 \\
        c - 2d = 0
    \end{cases}
    $$
    Let $b = s$ and $d = t$, then, $a = -2t, c = 2t$,
    $$E(A, 1) = \left\{\bmat{-2t \\ s \\ 2t \\ t} \suchthat s, t \in \RR\right\} = \spn\left\{\bmat{0 \\ 1 \\ 0 \\ 0}, \bmat{-2 \\ 0 \\ 2 \\ 1}\right\}$$
    The two elements are not scalar multiples of each other. Hence, $\left\{\bmat{0 \\ 1 \\ 0 \\ 0}, \bmat{-2 \\ 0 \\ 2 \\ 1}\right\}$ is a basis for $E(A, 1)$. With $a(1) = 2 = g(1)$.

    \item $E(A, 2)$
    $$\begin{amatrix}{4}
        -1 & 0 & 0 & 0 & 0 \\ 
        0 & -1 & 5 & -10 & 0 \\
        1 & 0 & 0 & 0 & 0 \\
        1 & 0 & 0 & 0 & 1
    \end{amatrix} \xrightarrow{\text{RREF}}
    \begin{amatrix}{4}
        1 & 0 & 0 & 1 & 0 \\
        0 & 1 & -5 & 0 & 0 \\
        0 & 0 & 0 & 1 & 0 \\
        0 & 0 & 0 & 0 & 0 
    \end{amatrix} \longleftrightarrow
    \begin{cases}
        a + d = 0 \\
        b - 5c = 0 \\
        d = 0
    \end{cases}
    $$
    Let $c = t$, then, $a = 0, b = 5t$,
    $$E(A, 2) = \left\{\bmat{0 \\ 5t \\ t \\ 0} \suchthat t \in \RR\right\} = \spn\left\{\bmat{0 \\ 5 \\ 1 \\ 0}\right\}$$
    And since there is only one nonzero vector in the set it is a basis for $E(A, 2)$. And $a(2) = 1 = g(2)$.

    \item $E(A, 3)$
    $$\begin{amatrix}{4}
        -2 & 0 & 0 & 0 & 0 \\
        0 & -2 & 5 & -10 & 0 \\
        1 & 0 & -1 & 0 & 0 \\
        1 & 0 & 0 & 0 & 0 
    \end{amatrix} \xrightarrow{\text{RREF}}
    \begin{amatrix}{4}
        1 & 0 & 0 & 0 & 0 \\
        0 & 1 & 0 & 5 & 0 \\
        0 & 0 & 1 & 0 & 0 \\
        0 & 0 & 0 & 0 & 0
    \end{amatrix} \longleftrightarrow
    \begin{cases}
        a = 0 \\
        b + 5d = 0 \\
        c = 0
    \end{cases}
    $$
    Let $d = t$, then, $b = -5t$,
    $$E(A, 3) = \left\{\bmat{0 \\ -5t \\ 0 \\ t} \suchthat t \in \RR\right\} = \spn\left\{\bmat{0 \\ -5 \\ 0 \\ 1}\right\}$$
    And since there is only one nonzero element in the set, it is a basis for $E(A, 3)$. And $a(3) = 1= g(3)$.

    From this let us form the set $E$ which contains the elements of the basis of each eigenspace $E(A, \lambda_i)$ for all $\lambda_i$.\footnote{One could check that set $E$ is linearly independent.}
    $$E = \left\{\bmat{0 \\ 1 \\ 0 \\ 0}, \bmat{-2 \\ 0 \\ 2 \\ 1}, \bmat{0 \\ 5 \\ 0 \\ 1}, \bmat{0 \\ -5 \\ 0 \\ 1}\right\}$$
    And since $|E| = 4$, which is equal to the size of $A$. Thus, $A$ is diagonalizable with $D = P\inv AP$ where,
    $$D = \bmat{1 & 0 & 0 & 0 \\ 0 & 1 & 0 & 0 \\ 0 & 0 & 2 & 0 \\ 0 & 0 & 0 & 3}, \quad P = \bmat{0 & -2 & 0 & 0 \\ 1 & 0 & 5 & -5 \\ 0 & 2 & 1 & 0 \\ 0 & 1 & 0 & 1}$$
    
\end{itemize}
\newpage
The permutations of columns in $P$ must follow the placing of eigenvalues in $D$. Any permutation for muliplicities (i.e. $\sigma(A) = \{1, 2, 1, 3\} \Longrightarrow D = \diag(1, 2, 1, 3)$) are allowed. To further illustrate this let's find another solution for the example above.

$$D = \diag(1, 2, 1, 3) = \bmat{1 & 0 & 0 & 0 \\ 0 & 2 & 0 & 0 \\ 0 & 0 & 1 & 0 \\ 0 & 0 & 0 & 3}, \quad P = \bmat{0 & 0 & -2 & 0 \\ 1 & 5 & 0 & -5 \\ 0 & 1 & 2 & 0 \\ 0 & 0 & 1 & 1}$$

One could check that this is true,\footnote{This can be easily verified using Matlab or Numpy.} and any permutation of the diagonal entries as long as the placing of vectors from $E$ must follow the order in $D$.



\begin{theorem}
Let $A \in M_n(\RR)$. Then $A$ is diagonalizable iff both the following conditions hold:
\begin{enumerate}
    \item Every root of $p_A(x)$ is real.
    \item For every eigenvalue $\lambda$ of $A$, $a(\lambda) = g(\lambda)$.
\end{enumerate}
\end{theorem}

\begin{remark}
Let $A \in M_n(\RR)$.
\begin{enumerate}
    \item For every eigenvalue $\lambda$ of $A_i$, $a(\lambda) \geq g(\lambda) \geq 1$.
    \item If $p_A(x)$ has exactly $n$ distinct real roots, then $A$ is diagonalizable.
\end{enumerate}
\end{remark}

\lline

\ul{\textbf{Recall.}}
\begin{enumerate}
    \item Standard inner product on $M_{n \times 1}(\RR)$,
    $$\langle x, y \rangle = x^{\mathsf{T}} y, \quad \forall x, y \in M_{n \times 1}(\RR)$$
    \item $S = \{v_1, v_2, \dots, v_n\}$ is an orthonormal set iff $\langle v_i, v_j \rangle = \begin{cases}
        0, \quad i \neq j \\
        1, \quad i = j
    \end{cases} = \delta_{ij}$
\end{enumerate}

\ul{\textbf{Question.}} When are the columns of $P \in M_n(\RR)$ orthonormal under the standard inner product of $M_{n \times 1}(\RR)$?

Let $P \in M_n(\RR)$, with columns $P_1, P_2, \dots, P_n$. Then $\{P_1, P_2, \dots, P_n\}$ is an orthonormal set iff:
$$\langle P_i, P_j \rangle = \delta_{ij}, \quad \forall i, j$$
$$P_i^{\mathsf{T}} P_j = \delta_{ij}, \quad \forall i, j$$
$$(\text{row } i \text{ of } P^{\mathsf{T}})(\text{column } j \text{ of } P) = \delta_{ij},\quad \forall i, j$$
$$P^{\mathsf{T}} P = I_n$$

\begin{definition}
$P \in M_n(\RR)$ is \textbf{orthogonal} if $P^{\mathsf{T}} P = I_n$.
\end{definition}

\begin{theorem}
Let $P \in M_n(\RR)$. Then the following are equivalent:
\begin{enumerate}
    \item $P$ is orthogonal.
    \item $P$ is nonsingular with $P\inv = P^{\mathsf{T}}$.
    \item The columns of $P$ form an orthonormal basis for $M_{n \times 1}(\RR)$ (under the standard inner product).
\end{enumerate}
\end{theorem}

\newpage
\begin{example}
Shown below are some examples of orthogonal matrices that can be verified by simply pairing the rows and getting their inner product (respectively pairing the columns and getting their inner product). It is not hard to show that the inner product is 0.
\leavevmode
\begin{enumerate}
    \item $$\bmat{\cos\theta & -\sin\theta \\ \sin\theta & \cos\theta}, \quad \forall \theta \in \RR$$
    \item $$\bmat{-2/3& 2/3 & 1/3 \\ 2/3 & 1/3 & 2/3 \\ 1/3 & 2/3 & -2/3}$$
\end{enumerate}
\end{example}

\begin{remark}
Let $P, Q \in M_n(\RR)$ be orthogonal
\leavevmode
\begin{enumerate}
    \item Then $PQ$ is also orthogonal
    $$(PQ)^{\mathsf{T}}(PQ) = Q^{\mathsf{T}} P^{\mathsf{T}} PQ = I_n$$
    \item $\det P = \pm 1$
    \begin{align*}
    (\det P^{\mathsf{T}})(\det P) &= \det I_n \\
    (\det P)^2 &= 1 \\
    \therefore \det P &= \pm 1
    \end{align*}
    \item Let $D \in M_n(\RR)$ be diagonal and $A = PDP^{\mathsf{T}}$. Then
    \begin{align*}
    A^{\mathsf{T}} &= (PDP^{\mathsf{T}})^{\mathsf{T}} \\
    &=(P^{\mathsf{T}})^{\mathsf{T}} D^{\mathsf{T}} P^{\mathsf{T}} \\
    &=PDP\inv = A
    \end{align*}
\end{enumerate}
\end{remark}

\begin{definition}\label{def:14.4}
Let $A \in M_n(\RR)$. Then
\begin{enumerate}
    \item $A$ is \textbf{symmetric} if $A^{\mathsf{T}} = A$.
    \item $A$ is \textbf{orthogonally diagonalizable} if there exists an orthogonal $P \in M_n(\RR)$ such that $P^{\mathsf{T}} AP$ is diagonal.
\end{enumerate}
\end{definition}

\begin{theorem}[Spectral Decomposition Theorem for Symmetric Matrices]
\leavevmode
$A \in M_n(\RR)$ is orthogonally diagonalizable iff $A$ is symmetric.
\end{theorem}

\lline

To orthogonally diagonalize a symmetric matrix $A \in M_n(\RR)$.
\begin{enumerate}
    \item Find the eigenvalues $\lambda_i$ of $A$.
    \item Find \textbf{orthonormal basis} of $E(A, \lambda_i)$ using Gram-Schmidt Process.
    \item Form the union $E = \{v_1, v_2, \dots, v_n\}$ of the orthonormal bases obtained in (2). Then $E$ is also an orthonormal basis for $M_{n\times 1}(\RR)$.
    \item $D = P^{\mathsf{T}} AP$ where $D = \diag(\lambda_1, \lambda_2, \dots, \lambda_n)$ and $P = \bmat{v_1 & v_2 & \cdots & v_n}$.
\end{enumerate}


\begin{example}
Let $A = \bmat{0 & 2 & 2 \\ 2 & 0 & 2 \\ 2 & 2 & 0}$. Find an orthogonal matrix $P$ and diagonal matrix $D$ such that $D = P^{\mathsf{T}} AP$.

\textit{Solution.}

Find the characteristic polynomial $p_A(x)$.
\item $$p_A(x) = \vmat{-x & 2 & 2 \\ 2 & -x & 2 \\ 2 & 2 & -x} = -(x+2)^2(x - 4)$$
Therefore, the eigenvalues of $A$: $-2, 4$

\begin{itemize}
    \item $E(A, -2)$
    $$\begin{amatrix}{3}
        2 & 2 & 2 & 0 \\
        2 & 2 & 2 & 0 \\
        2 & 2 & 2 & 0 
    \end{amatrix} \xrightarrow{\text{RREF}}
    \begin{amatrix}{3}
        1 & 1 & 1 & 0 \\ 
        0 & 0 & 0 & 0 \\
        0 & 0 & 0 & 0 \\
        0 & 0 & 0 & 0
    \end{amatrix} \longleftrightarrow
    \begin{cases}
        a + b + c = 0 \\
    \end{cases}
    $$
    Let $b = s$ and $c = t$, then, $a = -s - t$,
    $$E(A, -2) = \left\{\bmat{-s-t \\ s \\ t} \suchthat s, t \in \RR\right\} = \spn\left\{\bmat{-1 \\ 1 \\ 0}, \bmat{-1 \\ 0 \\ 1}\right\}$$
    Because they are not scalar multiples of each other, it is a basis for $E(A, -2)$.

    \begin{enumerate}
        \item Put $w_1 = v_1$,
        $$w_1 = \bmat{-1 \\ 1 \\ 0}$$
        And since $\lVert w_1 \rVert = \sqrt{2}$,
        $$w_1' = \frac{1}{\sqrt{2}}\bmat{-1 \\ 1 \\ 0} = \bmat{-1/\sqrt{2} \\ 1/\sqrt{2} \\ 0}$$

        \item Put $w_2 = v_2 - \dfrac{\langle v_2, w_1\rangle}{\langle w_1, w_1 \rangle}w_1$,
        $$w_2 = \bmat{-1 \\ 0 \\ 1} - \frac{1}{2}\bmat{-1 \\ 1 \\ 0} = \bmat{-1/2 \\ -1/2 \\ 1}$$
        And since $\lVert w_2 \rVert = \sqrt{\dfrac{3}{2}}$,
        $$w_2' = \sqrt{\frac{2}{3}}\bmat{-1/2 \\ -1/2 \\ 1} = \bmat{-1/\sqrt{6} \\ -1/\sqrt{6} \\ 2/\sqrt{6}}$$
    \end{enumerate}

    \item $E(A, 4)$
    $$\begin{amatrix}{3}
        -4 & 2 & 2 & 0 \\
        2 & -4 & 2 & 0 \\
        2 & 2 & -4 & 0 
    \end{amatrix} \xrightarrow{\text{RREF}}
    \begin{amatrix}{3}
        1 & 0 & -1 & 0 \\
        0 & 1 & -1 & 0 \\
        0 & 0 & 0 & 0
    \end{amatrix} \longleftrightarrow
    \begin{cases}
        a - c = 0 \\
        b - c = 0
    \end{cases}
    $$
    Let $c = t$, then, $a = t, b = t$,
    $$E(A, 4) = \left\{\bmat{t \\ t \\ t} \suchthat t \in \RR\right\} = \spn\left\{\bmat{1 \\ 1 \\ 1}\right\}$$
    And since it contains a nonzero element it is a basis for $E(A, 4)$. However, we should still normalize this. Since $\lVert w_1 \rVert = \sqrt{3}$, then,
    $$w_3' = \frac{1}{\sqrt{3}}\bmat{1 \\ 1 \\ 1} = \bmat{1/\sqrt{3} \\ 1/\sqrt{3} \\ 1/\sqrt{3}}$$
    Therefore, 
    $$\boxed{D = \bmat{-2 & 0 & 0 \\ 0 & -2 & 0 \\ 0 & 0 & 4}, \quad P = \bmat{-1/\sqrt{2} & -1/\sqrt{6} & 1/\sqrt{3} \\ 1/\sqrt{2} & -1/\sqrt{6} & 1/\sqrt{3} \\ 0 & 2/\sqrt{6} & 1/\sqrt{3}}}$$
\end{itemize}

\end{example}

\end{document}